% concatenated from SHARP_NHIM.tex
\documentclass{amsart}
\usepackage[utf8]{inputenc}
\usepackage{amsfonts, amssymb}
\usepackage{amssymb,amsmath,amscd,euscript,verbatim,array}
\usepackage{geometry}
\usepackage{anysize}
\usepackage{fancyhdr}
\usepackage{indentfirst}
\usepackage{graphicx}
\usepackage{color}
\usepackage{ifpdf}
\usepackage{marginnote}
\usepackage{enumitem}
\usepackage{mathrsfs}
\usepackage{amsthm}
%\usepackage{hyperref}
\usepackage{CJK}
\usepackage{bm}
\geometry{margin=1in}
\setlist[itemize]{leftmargin=2em}

%% Symbol Definitions --------------------------------------------------------------------



%% Theorem Definitions --------------------------------------------------------------------
%\newcounter{item}[section]
%\newcommand{\thenumber}{\stepcounter{item}\arabic{section}.\arabic{item}}
%\newcommand{\Proof}{\textbf{Proof}\hspace{0.3cm}}
%\newcommand{\End}{\ensuremath{\hfill{\Box}}\\}








%%Kaloshin's Definitions------------------------------------------------------------------
\newcommand{\N}{\mathbb{N}}
\newcommand{\Z}{\mathbb{Z}}
\newcommand{\R}{\mathbb{R}}
\newcommand{\Q}{\mathbb{Q}}
\newcommand{\C}{\mathbb{C}}
\newcommand{\Lip}{\mathrm{Lip}}
\newcommand{\Id}{\mathrm{Id}}
\newcommand{\ka}{\kappa}
\newcommand{\la}{\lambda}
\newcommand{\al}{\alpha}
\newcommand{\til}{\tilde}
\newcommand{\pa}{\partial}
\renewcommand{\th}{\theta}
\newcommand{\gm}{\gamma}

\newcommand{\Gm}{\Gamma}
\newcommand{\Th}{\Theta}
\newcommand{\A}{\mathbb{A}}
\newcommand{\cA}{\mathcal{A}}
\newcommand{\cM}{\mathcal{M}}
\newcommand{\cT}{\mathcal{T}}
\newcommand{\om}{\omega}
\newcommand{\sg}{\sigma}
\newcommand{\eps}{\varepsilon}
\newcommand{\dt}{\delta}
\newcommand{\lb}{\lambda}
\newcommand{\Lb}{\Lambda}
\newcommand{\cP}{\mathcal P}
\newcommand{\T}{\mathbb{T}}
\newcommand{\disp}{\displaystyle}
\newcommand{\D}[2]{\displaystyle \frac{\partial #1}{\partial #2}}
\newcommand{\Ds}[2]{\frac{\partial #1}{\partial #2}}
\newcommand{\DD}[3]{\displaystyle \frac{\partial^2 #1}{\partial #2\partial #3}}
\newcommand{\DDs}[3]{\frac{\partial^2 #1}{\partial #2\partial #3}}
\newcommand{\DDd}[2]{\displaystyle \frac{\partial^2 #1}{\partial {#2}^2}}
\newcommand{\DDds}[2]{\frac{\partial^2 #1}{\partial {#2}^2}}
\newcommand{\J}{\mathcal{J}}
\newcommand{\twist}{{\mathrm{twist}}}
\newcommand{\low}{{\mathrm{low}}}
\newcommand{\high}{{\mathrm{high}}}
\newcommand{\AM}{{\mathrm{AM}}}
\newcommand{\ovsp}{\mathrm{ovsp}}
\newcommand{\Ss}{\mathbb{S}}

\DeclareMathOperator{\supp}{supp}
\DeclareMathOperator{\osc}{osc}
\DeclareMathOperator{\diag}{diag}

%% Theorem Definitions --------------------------------------------------------------------




\newtheorem{Theorem}{Theorem}


\newtheorem{Definition}{Definition}[section]
\newtheorem{Proposition}[Definition]{Proposition}
\newtheorem{Corollary}[Theorem]{Corollary}
\newtheorem{Lemma}[Definition]{Lemma}
%\newtheorem{Figure}[Definition]{Figure}

%\theoremstyle{definition}
\newtheorem{Remark}[Theorem]{Remark}




%%Kaloshin's Definitions------------------------------------------------------------------


\newcommand{\bthm}{\begin{Theorem}}
	\newcommand{\ethm}{\end{Theorem}}
\newcommand{\bpr}{\begin{Proposition}}
	\newcommand{\epr}{\end{Proposition}}
\newcommand{\blm}{\begin{Lemma}}
	\newcommand{\elm}{\end{Lemma}}
\newcommand{\bex}{\begin{Exercise}}
	\newcommand{\eex}{\end{Exercise}}
\newcommand{\be}{\begin{equation}}
	\newcommand{\ee}{\end{equation}}
\newcommand{\beal}{\begin{aligned}}
	\newcommand{\enal}{\end{aligned}}
\newcommand{\brm}{\begin{Remark}}
	\newcommand{\erm}{\end{Remark}}




\begin{document}
	
	
	\title[On the sharpness of the \(C^1\)-norm threshold for perturbations]{On the sharpness of the \(C^1\)-norm threshold for perturbations
		in the normally hyperbolic invariant manifold theorem\\---a toy model perspective}
	
	\author{Yujie Huang}
	\address{School of Mathematics and Statistics, Beijing Institute of Technology, Beijing 100081, China}
	\email{huangyujie@bit.edu.cn}
	
	\author{Junhao Li}
	\address{School of Mathematics and Statistics, Beijing Institute of Technology, Beijing 100081, China}
	\email{junhaoli@bit.edu.cn}
	
	\author{Lin Wang}
	\address{School of Mathematics and Statistics, Beijing Institute of Technology, Beijing 100081, China}
	\email{lwang@bit.edu.cn}
	
	
	\subjclass[2010]{Primary 37J40; Secondary 37E40}
	
	\keywords{Dissipative twist maps, Normally hyperbolic invariant manifolds, Invariant graphs}
	
	
	
	\begin{abstract}
		
		The classical normally hyperbolic invariant manifold theorem asserts that a \(C^1\) normally hyperbolic invariant manifold persists under \(C^1\) small perturbations. For a family of standard-like dissipative twist maps,  we show that the threshold \((1-\sqrt{\lambda})^2\) for the \(C^1\)-norm of the perturbation is sharp: there exists a $C^\infty$ perturbation \(\phi\) with \(\|\phi\|_{C^1} = (1-\sqrt{\lambda})^2\) such that the map preserves a unique invariant graph, but this graph possesses  non-differentiable points. On the other hand, whenever \(\|\phi\|_{C^1} < (1-\sqrt{\lambda})^2\), the \(C^1\) normally hyperbolic invariant manifold persists, where \(\lambda\) denotes the Jacobian determinant of the map. This provides a critical threshold phenomenon for the persistence of invariant graphs in dissipative twist maps.
	\end{abstract}
	
	
	
	\maketitle
	
	\vspace{2em}
	
	%\tableofcontents
	%
	%	\newpage
	\section{\sc Introduction}
	Let $\alpha:=(\alpha_1,\alpha_2)\in \R^2$. Fix $\lambda\in (0,1)$. We consider
	\begin{equation}\label{F1}
		T_{\lambda,\alpha}(x,y):=(x+\alpha_1+\lambda y,\alpha_2+\lambda y).
	\end{equation}
	The parameter $\lambda$ controls the dissipation, and $\alpha_1,\alpha_2$
	are constants that determine the translation in the $x$ and $y$ directions, respectively. For any parameter pair $(\alpha_1,\alpha_2) \in \mathbb{R}^2$, a direct calculation shows that the integrable system \eqref{F1} possesses a unique invariant graph
	\begin{equation}\label{ingm1}
		\Gamma_0 = \mathbb{R} \times \left\{\frac{\alpha_2}{1-\lambda}\right\}.
	\end{equation}
	
	
	
	The map $T_{\la,\alpha}$ restricted on $\Gamma_0$ reduces to a circle diffeomorphism $g_\al$ with rotation number $\al_1+\frac{\lambda}{1-\lambda}\al_2$. Note that the invariant graph $\Gamma_0$  is normally hyperbolic. More precisely, it is  so called immediately absolutely $r$-normally hyperbolic for any $r\in\N$ in terms of \cite[Definition 2]{HPS}. This notion is also referred to as $r$-bunching (see for instance \cite{KH}).
	
	
	
	\begin{Definition}
		Let $V$ be a smooth compact submanifold of a smooth manifold $\mathcal M$.
		Suppose $f: \mathcal M \to \mathcal M$ is a $C^{1}$ diffeomorphism and $f(V) = V$.
		Let $T_V \mathcal M$, the tangent bundle of $\mathcal M$ over $V$, have a $Tf$-continuous invariant splitting
		\[
		T_V \mathcal M = N^u \oplus TV \oplus N^s .
		\]
		For any $z \in \mathcal M$ put
		\[
		T_z f = N^u_z f \oplus V_z f \oplus N^s_z f .
		\]
		Thus,
		\[
		Tf|_{TV} = Vf,\quad
		Tf|_{N^u} = N^u f,\quad
		Tf|_{N^s} = N^s f .
		\]
		$f$ is \emph{immediately absolutely $r$-normally hyperbolic at $V$} iff
		$f$ is $C^r$ and there is a Riemann structure on $T\mathcal M$ such that for all $z \in V$,
		$k=0$ and $r$:
		\begin{enumerate}
			\item[(a)] $\inf_z m(N^u_z f) > \sup_z \|V_z f\|^k$ \quad and
			\item[(b)] $\sup_z \|N^s_z f\| < \inf_z m(V_z f)^k$ ,
		\end{enumerate}
		where the minimum norm $m(A)$ of a linear transformation $A$ is defined as
		\[
		m(A) = \inf \{|Ax| : |x| = 1\} .
		\]
		When $A$ is invertible, $m(A) = \|A^{-1}\|^{-1}$.
	\end{Definition}
	
	
	For each $(x,y)\in\R^2$,
	\[DT_{\la,\alpha}(x,y)=\left(
	\begin{array}{cc}
		1 & \la \\
		0 & \la \\
	\end{array}
	\right)=P\left(
	\begin{array}{cc}
		1 & 0 \\
		0 & \la \\
	\end{array}
	\right)P^{-1},
	\]
	where $P=\left(
	\begin{array}{cc}
		1 & \frac{-\la}{1-\la} \\
		0 & 1 \\
	\end{array}
	\right)$. Hence, we have a $TT_{\la,\al}$-invariant splitting
	\[T\Gamma\oplus N^s,\]
	where $TT_{\la,\al}$ denotes the tangent map of $T_{\la,\al}$.
	Moreover, there exists a constant $C$ such that for all $i\in\Z$ and $z\in \Gamma$, there hold
	\begin{equation}\label{hype}
		\frac{1}{C}\la^i|v|\leq |DT_{\la,\al}^i(z)v|\leq C\la^i|v|,\quad \mathrm{for}\ v\in N^s_z,
	\end{equation}
	\[\frac{1}{C}|v|\leq |DT^i_{\la,\al}(z)v|\leq C|v|,\quad \mathrm{for}\ v\in T_z\Gamma,\]
	where the constant $C$ can be estimated by the norms of $P$ and $P^{-1}$.
	
	We consider the persistence of invariant graphs under perturbations. More precisely, we study the three-parameter family
	\begin{equation}\label{Fp}
		F^\phi_{\lambda,\alpha}(x,y) = \big(x + \alpha_1 + \lambda y + \phi(x),\; \alpha_2 + \lambda y + \phi(x)\big),
	\end{equation}
	where \(\lambda\in(0,1)\), \(\alpha=(\alpha_1,\alpha_2)\in\mathbb{R}^2\), and \(\phi\) is a \(1\)-periodic \(C^1\) function satisfying \(\int_{0}^{1} \phi(x)\,dx = 0\).  The map \(F^\phi\colon\mathbb{R}^2\to\mathbb{R}^2\) is the lift to the universal cover of a dissipative twist map on \(\mathbb{T}\times\mathbb{R}\). If \(\alpha_2=0\), then \(F^\phi\) is called an exact conformally symplectic twist map. For related studies on dissipative twist maps and exact conformally symplectic twist maps, we refer the reader to \cite{Cas87,CCD13,CCD20,CCD22,Lec87,Mas23} and the references therein.
	
	The classical normally hyperbolic invariant manifold (NHIM) theorem asserts that for fixed \(\lambda\) and any \(\alpha_1,\alpha_2\in\mathbb{R}\), if \(\|\phi\|_{C^1}\) is sufficiently small, then \(F^\phi\) preserves a \(C^1\) invariant graph, where
	\[\|\phi\|_{C^1}:=\max\{\|\phi\|_{C^0},\|\phi'\|_{C^0}\}.\]
	The NHIM theorem provides a sufficient condition for the existence of a \(C^1\) invariant graph. In this paper, we investigate the necessity aspect from a quantitative viewpoint, i.e., we address the following question:
	
	\vspace{1ex}
	
	\begin{itemize}
		\item \textbf{Question:} {\it For any \(\lambda\in(0,1)\), what is  the critical size of the \(C^1\)-norm of the perturbation \(\phi\) that guarantees a \(C^1\) invariant graph?}
	\end{itemize}
	
	\vspace{1ex}
	\subsection{Main results}
	Birkhoff's graph theorem \cite{Bir20} states that if \(F_{\lambda,\alpha}^\phi\) admits an invariant graph, then the graph must be at least Lipschitz. Moreover, it is easy to show that the invariant graph, if it exists, is necessarily unique. On the one hand, we give a quantitative version of the NHIM theorem for the model \eqref{Fp}.
	
	\begin{Theorem}\label{M2}
		Given \(\lambda \in (0,1)\), for any \(\alpha \in \mathbb{R}^2\) and any \(C^1\) perturbation \(\phi\) with zero average satisfying
		\begin{equation}\label{contrac}
			\|\phi\|_{\mathrm{Lip}} \leq (1-\sqrt{\lambda})^2,
		\end{equation}
		the map \(F_{\lambda,\alpha}^\phi\) admits a unique Lipschitz invariant graph. Moreover, if
		\begin{equation}\label{m22con}
			\|\phi\|_{C^1} < (1-\sqrt{\lambda})^2,
		\end{equation}
		then the map \(F_{\lambda,\alpha}^\phi\) admits a unique \(C^1\) invariant graph, and this invariant graph is $1$-normally hyperbolic.
	\end{Theorem}
	
	The proof of Theorem~\ref{M2} is rather standard. First, under condition \eqref{contrac}, the existence of a Lipschitz invariant graph is obtained via the classical graph transform method. Then, under condition \eqref{m22con}, one obtains that \(F_{\lambda,\alpha}^\phi\) is a \(1\)-fiber contraction (see the definition in \cite[Page 31]{HPS}). On this basis, the regularity of the invariant graph is upgraded to \(C^1\). This upgrade can be carried out using the Lipschitz jet technique (see the proof of the \(C^r\) section theorem in \cite[Pages 33--35]{HPS}), or by exploiting the geometric structure of the invariant graph itself (see \cite{BB}). Here, to better take advantage of the specific features of the model under consideration, we adopt the latter approach.
	
	On the other hand, we establish the following result.
	
	\vspace{1ex}
	
	\begin{Theorem}\label{Th1}
		Given \(\lambda \in (0,1)\), for any rational number \(\beta\), there exist \(\alpha \in \mathbb{R}^2\) and \(\phi_\lambda \in C^\infty(\mathbb{T})\) such that
		\[
		\|\phi_{\lambda}\|_{C^1} = (1-\sqrt{\lambda})^2,
		\]
		and \(F_{\lambda,\alpha}^{\phi_{\lambda}}\) admits a unique Lipschitz invariant graph. Moreover, the restriction of \(F_{\lambda,\alpha}^{\phi_{\lambda}}\) to this invariant graph is a circle map of frequency \(\beta\), but the invariant graph possesses non-differentiable points.
	\end{Theorem}
	The proof of Theorem~\ref{Th1} is constructive. The construction is based on the theory of circle maps \cite{H79} and the dissipative versions of the Herman-Mather formula and its derivation (see \eqref{eq:main} and (A') below), and draws on ideas from the constructions of Herman \cite{H83} and Arnaud \cite{Ar11} in the area-preserving twist map setting. The difference here is that we require not only that the invariant graph possesses non-differentiable points, but also that the \(C^1\)-norm of the perturbation has a prescribed a priori magnitude. Compared with the conservative case, this imposes additional constraints.
	
	Combining Theorem \ref{M2} and Theorem~\ref{Th1}, one knows that the threshold \((1-\sqrt{\lambda})^2\) for the \(C^1\)-norm of the perturbation is sharp: there exists a \(C^\infty\) perturbation \(\phi_{\lambda}\) with \(\|\phi_{\lambda}\|_{C^1} = (1-\sqrt{\lambda})^2\) such that the map preserves a unique invariant graph, yet this graph possesses non-differentiable points. On the other hand, whenever \(\|\phi\|_{C^1} < (1-\sqrt{\lambda})^2\), the \(C^1\) normally hyperbolic invariant manifold persists.
	
	\begin{Remark}\label{R2}
		Fix \(\lambda \in (0,1)\). We define the following sets:
		\[
		\begin{aligned}
			A &:= \bigl\{ \phi \in C^1(\mathbb{T}) \;\big|\; \phi \text{ satisfies } \eqref{m22con} \bigr\}, \\
			B &:= \bigl\{ \phi \in C^1(\mathbb{T}) \;\big|\; \exists\,  \alpha \in \mathbb{R}^2 \text{ s.t. the invariant graph } \Gamma \text{ of } F_{\lambda,\alpha}^\phi \text{ is } 1\text{-normally hyperbolic} \bigr\}, \\
			C &:= \bigl\{ \phi \in C^1(\mathbb{T}) \;\big|\; \exists\, \alpha \in \mathbb{R}^2 \text{ s.t. } F_{\lambda,\alpha}^\phi \text{ admits a } C^1 \text{ invariant graph } \Gamma \bigr\}.
		\end{aligned}
		\]
		The relationship among the sets \(A, B, C\) and the perturbation \(\phi\) constructed in Theorem~\ref{Th1} is illustrated in Fig.~\ref{X}. For the detailed verification, see Section~\ref{R3}.
		
		\begin{figure}[htbp]
			\small \centering
			\includegraphics[width=5.5cm]{X.png}
			\caption{Relationship among the sets \(A, B, C\)}
			\label{X}
		\end{figure}
		
		The region enclosed by the orange curve represents the set \(A\), the region enclosed by the blue curve represents the set \(B\), and the region enclosed by the green curve represents the set \(C\). The red points on the boundary of \(A\) correspond to the perturbations \(\phi_\lambda\) constructed in Theorem~\ref{Th1}. Note that Theorem~\ref{Th1} shows that such red points are at least countably infinite on the boundary of \(A\).
	\end{Remark}
	
	
	\begin{Remark}
		In \cite[Page 1]{Man78}, Ma\~{n}\'{e} introduced the ``persistent'' condition. To give a concrete definition of the ``persistent'' condition in our setting, we pass from the universal covering space back to the cylinder. Let \(\mathcal M = \mathbb{T} \times \mathbb{R}\), and let \(\mathrm{Diff}^1(\mathcal M)\) denote the set of diffeomorphisms from \(\mathcal M\) to itself. Denote by \(\gamma\) and \(f\) the projections of \(\Gamma\) and \(F_{\lambda,\alpha}^\phi\) onto \(\mathcal M\), respectively. The invariant graph \(\gamma\) is called {persistent} if there exist a neighborhood \(\mathcal{U}\) of \(\gamma\) in \(\mathcal M\) and a neighborhood \(\mathcal{V}\) of \(f\) in \(\mathrm{Diff}^1(\mathcal M)\) such that
		\begin{enumerate}
			\item[(a)] for every \(\tilde{f} \in \mathcal{V}\), the set \(\gamma_{\tilde{f}} = \bigcap_{n \in \mathbb{Z}} \tilde{f}^n(\mathcal{U})\) is a \(C^1\) submanifold of \(\mathcal M\) and \(\gamma_f = \gamma\);
			\item[(b)] \(\gamma_{\tilde{f}}\) is \(C^1\) close to \(\gamma\) whenever \(\tilde{f}\) is \(C^1\) close to \(f\).
		\end{enumerate}
		In our setting, we say that \(\Gamma\) is persistent if its projection \(\gamma\) is persistent. Moreover, we may define
		\[
		B' := \bigl\{ \phi \in C^1(\mathbb{T}) \;\big|\; \exists\, \alpha \in \mathbb{R}^2 \text{ such that the invariant graph } \Gamma \text{ of } F_{\lambda,\alpha}^\phi \text{ is persistent} \bigr\}.
		\]
		According to \cite{Man78} and the NHIM theorem, the set \(B'\) coincides with the set \(B\) defined in Remark~\ref{R2}.
	\end{Remark}
	
	
	\subsection{A key ingredient: the Herman-Mather formula}
	Following Herman \cite{H79}, we denote by \(\mathrm{Diff}^r_+(\mathbb{R})\) (resp. \(\mathrm{Diff}^r_+(\mathbb{T})\)) the group of orientation-preserving \(C^r\) diffeomorphisms on \(\mathbb{R}\) (resp. \(\mathbb{T}\)), where \(r \in [0,+\infty) \cup \{+\infty\} \cup \{\omega\}\). The universal covering space of \(\mathrm{Diff}^r_+(\mathbb{T})\) can be represented as
	\[
	D^r(\mathbb{T}) := \{f \in \mathrm{Diff}^r_+(\mathbb{R}) \mid f - \mathrm{Id} \in C^r(\mathbb{T})\}.
	\]
	For any \(f \in D^r(\mathbb{T})\), the rotation number \(\rho(f)\) is well-defined. By \cite[Proposition 2.1]{SoW},
	the map \(F_{\lambda,\alpha}^{\phi_{\lambda}}\) admits an invariant graph \(\Gamma := \{(x, \Psi_{\lambda,\alpha}(x)) \mid x \in \mathbb{R}\}\) if and only if there exists \(g_{\lambda,\alpha} \in D^r(\mathbb{T})\) ($r\geq 0$) satisfying the following equation
	\begin{equation}\tag{A}\label{eq:main}
		g_{\lambda,\alpha}(x) + \lambda g_{\lambda,\alpha}^{-1}(x) = (1+\lambda)x + (1-\lambda)\alpha_1 + \lambda\alpha_2 + \phi_{\lambda}(x) \quad \forall x \in \mathbb{R}.
	\end{equation}
	Note that
	\[
	g_{\lambda,\alpha}(x) = x + \alpha_1 + \lambda \Psi_{\lambda,\alpha}(x) + \phi_\lambda(x).
	\]
	For $\lambda = 1$, formula (A) was given by Herman in \cite[Chapter II]{H83}, based on which he conducted a detailed study of the dynamics of area-preserving twist maps. Meanwhile, Mather established the same result in \cite{M84} independently. Recently, \cite{SoW} extended this formula to the dissipative setting. Henceforth, we shall refer to formula (A) as the {\it Herman--Mather formula}.
	
	
	If \(g_{\lambda,\alpha}\) is differentiable at \(x\), differentiating both sides of \eqref{eq:main} with respect to \(x\) yields
	\begin{equation}\tag{A'}\label{gdiff}
		g'_{\lambda,\alpha}(x) + \frac{\lambda}{g'_{\lambda,\alpha}(g_{\lambda,\alpha}^{-1}(x))} = 1+\lambda + \phi_\lambda'(x).
	\end{equation}
	We shall call formula (A') the {\it derived Herman--Mather formula}. Essentially, the Herman--Mather formula is derived from the invariance of the graph; its significance lies in the fact that {\it it directly relates the perturbation of the map to the dynamics of the map restricted to the invariant graph, while the invariant graph itself does not appear explicitly in the formula.}
	
	Most of the proofs and constructions in this paper rely closely on this formula. In particular, it is used in the verification of the normal hyperbolicity of the invariant graph in the proof of Theorem 1 (Lemma~\ref{lem:universal-g-lower}), in the construction of perturbations in Theorem 2, and in the proof of the strict inclusion relations between the sets in Remark~\ref{R2}.
	
	
	\vspace{1ex}
	\noindent\textbf{Notation convention.} For notational simplicity, we shall omit the subscripts and superscripts of the relevant notations whenever no confusion arises. For instance, we write \(F\), \(\phi\), \(\Psi\), \(g\), etc., instead of the more cumbersome \(F_{\lambda,\alpha}^{\phi^\lambda}\), \(\phi^\lambda\), \(\Psi_{\lambda,\alpha}\), \(g_{\lambda,\alpha}\), and so on. The dependence of these quantities on \(\lambda\) and \(\alpha\) should be clear from the context.
	
	\section{\sc Proof of Theorem \ref{M2}}
	\subsection{Persistence of Lipschitz invariant graph}
	
	\subsubsection{Construction of the graph transform}
	Fix $\la\in (0,1)$. Let \begin{equation}\label{iI}
		I := \left[-\frac{\al_2}{1-\la}-(1-\lambda), \frac{\al_2}{1-\la}+(1-\lambda)\right]
	\end{equation}and denote by $\Lip(\T, I)$ the space of 1-periodic Lipschitz maps $\Psi: \R \to I$ with Lipschitz constant $\mathcal{L}(\Psi)$. For $K > 0$, define
	\[
	\Lip_K := \{\Psi \in \Lip(\T, I) \mid \mathcal{L}(\Psi) \leq K\}.
	\]
	
	Consider the map ${F_{\lambda,\alpha}^\phi}(x,y) = (X(x,y), Y(x,y))$ where
	\begin{align*}
		X(x,y) &:= x + \alpha_1 + \lambda y + \phi(x), \\
		Y(x,y) &:= \al_2+\lambda y + \phi(x).
	\end{align*}
	For any $\psi \in \Lip_K$, we have the compositions
	\begin{align*}
		X \circ (\mathrm{Id}, \Psi)(x) &= x + \alpha_1 + \lambda \Psi(x) + \phi(x), \\
		Y \circ (\mathrm{Id}, \Psi)(x) &= \al_2+\lambda \Psi(x) + \phi(x).
	\end{align*}
	Let $A_\lambda := \|\phi\|_{\mathrm{Lip}}$. It follows from $\int_\T\phi(x)dx=0$ that $\|\phi\|_{C^0}\leq A_\lambda$.
	
	\begin{Lemma}\label{invertib}
		Define the constant
		\[
		K_1 := \frac{2}{\sqrt{\lambda}} - 1.
		\]
		If $K < K_1$, then for each $\Psi \in \Lip_K$, the map $X \circ (\mathrm{Id}, \Psi)$ is invertible with Lipschitz inverse satisfying
		\[
		\mathcal{L}\left([X \circ (\mathrm{Id}, \Psi)]^{-1}\right) \leq \frac{1}{1 - \lambda K - A_\lambda}.
		\]
	\end{Lemma}
	
	\begin{proof}
		Define $\tilde{X}(x) := X \circ (\mathrm{Id}, \Psi)(x) - x=\alpha_1+\lambda\Psi(x)+\phi(x)$. Then $\tilde{X}$ is Lipschitz with
		\[
		|\tilde{X}(x_1) - \tilde{X}(x_2)| \leq (\lambda K + A_\lambda)|x_1 - x_2|.
		\]
		The condition $K < K_1$ implies $\lambda K + (1 - \sqrt{\lambda})^2 < 1$, and consequently
		\[
		\mathcal{L}(\tilde{X}) \leq \lambda|\Psi(x_1)-\Psi(x_2)|+|\phi(x_1)-\phi(x_2)|\leq \lambda K + A_\lambda < \lambda K + (1 - \sqrt{\lambda})^2 < 1.
		\]
		Therefore, the map $X\circ(\mathrm{Id},\Psi)$ is invertible and the Lipschitz constant of $[X\circ(\mathrm{Id},\Psi)]^{-1}$ satisfies:
		\[
		\mathcal{L}([X\circ(\mathrm{Id},\Psi)]^{-1})\leq \frac{1}{1-\lambda K -A_{\lambda}}
		\]
	\end{proof}
	
	When $[X \circ (\mathrm{Id}, \Psi)]^{-1}$ exists, we define the graph transform $\mathcal{T}: \Lip_K \to \Lip_K$ by
	\[
	\mathcal{T}\Psi : x \mapsto Y \circ (\mathrm{Id}, \Psi) \circ [X \circ (\mathrm{Id}, \Psi)]^{-1}(x).
	\]
	For the graph $\tilde{\Gamma}(\Psi) := \{(x, \Psi(x)) \mid x \in \R\}$, we have the invariance property
	\[
	\tilde{\Gamma}(\mathcal{T}\Psi) = {F_{\lambda,\alpha}^\phi}(\tilde{\Gamma}(\Psi)).
	\]
	
	\begin{Lemma}\label{welld}
		Define the constant
		\[
		K_2 := \frac{1}{\sqrt{\lambda}} - 1.
		\]
		The graph transform $\mathcal{T}: \Lip_{K_2} \to \Lip_{K_2}$ is well-defined, i.e., $\mathcal{T}\Psi \in \Lip_{K_2}$ for all $\Psi \in \Lip_{K_2}$.
	\end{Lemma}
	
	\begin{proof}
		For $\Psi \in \Lip_{K_2}$, we first establish the uniform bound:
		\[
		\|\mathcal{T}\Psi\|_{C^0} \leq \al_2+\lambda \|\Psi\|_{C^0} + A_\lambda \leq \al_2+\frac{\la\al_2}{1-\la}+\lambda(1-\lambda) + (1 - \sqrt{\lambda})^2 <\frac{\al_2}{1-\la}+ 1-\lambda.
		\]
		
		The Lipschitz estimate follows from:
		\begin{align*}
			|\mathcal{T}\Psi(x_1) - \mathcal{T}\Psi(x_2)|
			&\leq \frac{\lambda \mathcal{L}(\Psi) + A_\lambda}{1 - \lambda \mathcal{L}(\Psi) - A_\lambda}|x_1 - x_2| \\
			&< \frac{\lambda K_2 + (1 - \sqrt{\lambda})^2}{1 - \lambda K_2 - (1 - \sqrt{\lambda})^2}|x_1 - x_2| \\
			&= K_2 |x_1 - x_2|,
		\end{align*}
		the first inequality above is obtained from the definition of $\mathcal{T}$ and Lemma~\ref{invertib}.
		
		Thus $\mathcal{T}\Psi \in \Lip_{K_2}$.
	\end{proof}
	
	
	
	
	\subsubsection{Contraction mapping}
	Note that $\Lip_K$ is a closed subspace of the Banach space $C^0(\T,I)$ equipped with the $C^0$-metric, and hence is complete. To complete the proof of Theorem~\ref{M2}, it remains to show that the graph transform is a contraction mapping. Assuming the invertibility of $X \circ (\Id,\psi)$ and the well-definedness of $\mathcal{T}(\Lip_K\to \Lip_K)$, we establish this through the following lemma.
	
	\begin{Lemma}\label{ctrmap}
		Define the constant
		\[
		K_3 := \frac{1}{\lambda} - 1.
		\]
		If $K < K_3$, then the graph transform $\mathcal{T}: \Lip_K \to \Lip_K$ is a contraction. Specifically, for any $\Psi_1, \Psi_2 \in \Lip_K$,
		\[
		\|\mathcal{T}\Psi_1 - \mathcal{T}\Psi_2\|_{C^0} \leq l \|\Psi_1 - \Psi_2\|_{C^0},
		\]
		where $0 < l < 1$.
	\end{Lemma}
	
	\begin{proof}
		Fix $z \in \T$ and $\Psi_1, \Psi_2 \in \Lip_K$. Let $(x,y)$ be the point on the graph of $\Psi_1$ determined by
		\[
		x := [X \circ (\Id, \Psi_1)]^{-1}(z), \quad y := \Psi_1(x).
		\]
		
		By definition of the graph transform, we have:
		\begin{align*}
			&\mathcal{T}\Psi_1(z) = Y\circ(\Id, \Psi_1)(x), \\
			&\mathcal{T}\Psi_2(z) = \mathcal{T}\Psi_2 \circ X\circ(\Id, \Psi_1)(x),\\
			&\mathcal{T}\Psi_2 \circ X(x, \Psi_2(x)) =Y\circ(\Id, \Psi_2)(x).
		\end{align*}
		It follows that
		\[|Y\circ(\Id, \Psi_2)(x) - \mathcal{T}\Psi_2 \circ X\circ(\Id, \Psi_1)(x)|\leq K\lambda |\Psi_1(x) - \Psi_2(x)|.\]
		Moreover, we have
		\begin{align*}
			|\mathcal{T}\Psi_1(z) - \mathcal{T}\Psi_2(z)|
			&\leq |Y\circ(\Id, \Psi_1)(x) - Y\circ(\Id, \Psi_2)(x)| \\
			&\quad + |Y\circ(\Id, \Psi_2)(x) - \mathcal{T}\Psi_2 \circ X\circ(\Id, \Psi_1)(x)| \\
			&\leq \lambda |\Psi_1(x) - \Psi_2(x)| + K\lambda |\Psi_1(x) - \Psi_2(x)| \\
			&= (\lambda + K\lambda) |\Psi_1(x) - \Psi_2(x)|.
		\end{align*}
		This yields the uniform estimate:
		\[
		\|\mathcal{T}\Psi_1 - \mathcal{T}\Psi_2\|_{C^0} \leq l \|\Psi_1 - \Psi_2\|_{C^0},
		\]
		where $l := \lambda(1 + K)$. When $K < K_3 = \frac{1}{\lambda} - 1$, we have $l < 1$, proving that $\mathcal{T}$ is indeed a contraction.
	\end{proof}
	
	
	
	
	
	
	
	
	
	
	
	Note that for each $\lambda\in (0,1)$, we have
	\[K_2<\min\{K_1,K_3\}.\]
	It follows that the map $\mathcal{T} \colon \mathrm{Lip}_{K_2} \to \mathrm{Lip}_{K_2}$ admits a unique fixed point $\tilde{\Psi}$. Moreover, the graph
	\begin{equation}\label{tGp}
		\tilde{\Gamma}(\tilde{\Psi}) := \bigl\{ \bigl(x, \tilde \Psi(x)\bigr) \bigm| x \in \mathbb{R} \bigr\}
	\end{equation}
	is the unique Lipschitz invariant graph for ${F_{\lambda,\alpha}^\phi}$.
	
	
	
	
	
	
	\subsection{Persistence of $C^1$ invariant graph}
	The proof  is inspired by  \cite[Proof of Theorem 3.1]{BB}). We proceed as follows.
	\subsubsection{The cone condition}
	Denote
	\[\zeta:=\frac{1}{\sqrt{\lambda}}-1.\]
	We denote $F$ for ${F_{\lambda,\alpha}^\phi}$, $\Psi$ for $\tilde{\Psi}$, and we use $\Gamma$ instead of $\tilde{\Gamma}$ to denote the invariant graph by $F$  for simplicity.
	For $z\in\R^2$, we consider the cone  as follows:
	\[\mathcal{C}_\zeta(z):=\{v=(v_1,v_2)\in\R^2\ |\ |v_2|\leq \zeta |v_1|\},\]
	where we identified the Euclidean space $\R^2$ and its tangent space $\mathrm{T}_z\R^2$. The inner product in $\R^2$ is given by the standard one. Recall the projection $\pi_1:\R^2\to\R$ via
	\[\pi_1:(x,\Psi(x))\mapsto x,\]
	which means $\pi_1^{-1}(x)\in \Gamma$ for each $x\in\R$. Moreover, we denote the cone along $\Gamma$ by $\mathcal{C}_\zeta(x)$ instead of  $\mathcal{C}_\zeta(\pi_1^{-1}(x))$ for simplicity. Since the invariant graph $\Gamma$ is {\it a priori} only Lipschitz, we also need to consider the (Bouligand) tangent cone along $\Gamma=\{(x,\Psi(x))\ |\ x\in\R\}$.
	\[TC_\Gamma(x):=\left\{v\in\R^2\ \left|\ v=\mu\lim_{n\to+\infty}\frac{(x_n,\Psi(x_n))^T-(x,\Psi(x))^T}{\|(x_n,\Psi(x_n))^T-(x,\Psi(x))^T\|}\right.,\quad \forall \mu\in\R, x_n\to x \in \R\right\},\]
	where  $a^T$ denotes the transpose of $a$. Let us recall that there exists $g\in D^0(\T)$ such that
	\[F(x,\Psi(x))=(g(x),\Psi(g(x))).\]
	A direct calculation implies that for each $v\in TC_\Gamma(x)$, there exists $\mu\in\R$ such that
	\[DF(\pi_1^{-1}(x))v=\mu\lim_{n\to+\infty}\frac{(g(x_n),\Psi(g(x_n)))^T-(g(x),\Psi(g(x)))^T}{\|(g(x_n),\Psi(g(x_n)))^T-(g(x),\Psi(g(x)))^T\|}.\]
	Note that $\Gamma$ is invariant by $F$. Then its tangent cone is also invariant by the tangent map $DF$.
	It follows that
	\begin{equation}\label{cooncd}
		DF(\pi_1^{-1}(x))\cdot\left(TC_\Gamma(x)\right):=\left\{DF(\pi_1^{-1}(x))v\ |\ v\in TC_\Gamma(x)\right\}= TC_\Gamma(g(x)),
	\end{equation}
	
	Next, we need to verify the following two conditions:
	\begin{itemize}
		\item [(I)] $TC_\Gamma(x)\subseteq \mathcal{C}_\zeta(x)$ for each $x\in \R$;
		\item [(II)] $DF(\pi_1^{-1}(x))\cdot(\mathcal{C}_\zeta(x))\subseteq \mathrm{interior}\ \mathcal{C}_\zeta(x)\cup \{(0,0)\}$  for each $x\in \R$.
	\end{itemize}
	By Lemma \ref{welld} and (\ref{tGp}), we know that $\mathcal{L}(\Psi)\leq \frac{1}{\sqrt{\lambda}}-1$, which means Item (I) holds.
	For Item (II), it follows directly from the definition of $F$ that for any vector $v = (v_1, v_2) \in \mathcal{C}_\zeta(x)$, we have
	\[
	DF(\pi_1^{-1}(x))v =
	\begin{pmatrix}
		1+\phi'(x) & \lambda \\
		\phi'(x) & \lambda
	\end{pmatrix}
	\begin{pmatrix}
		v_1 \\
		v_2
	\end{pmatrix}
	=
	\begin{pmatrix}
		v_1 + \phi'(x)v_1 + \lambda v_2 \\
		\phi'(x)v_1 + \lambda v_2
	\end{pmatrix}.
	\]
	Note that $\|\phi\|_{C^1} < (1 - \sqrt{\lambda})^2$ and $|v_2 / v_1| \leq \zeta = \frac{1}{\sqrt{\lambda}} - 1$, so $|\phi'(x)+\lambda \frac{v_2}{v_1}|<1-\sqrt{\lambda}<1$. It is straightforward to verify that
	\begin{equation}\label{test2}
		\frac{|\phi'(x)v_1 + \lambda v_2|}{|v_1 + \phi'(x)v_1 + \lambda v_2|}
		= \frac{|\phi'(x) + \lambda \frac{v_2}{v_1}|}{|1 + \phi'(x) + \lambda \frac{v_2}{v_1}|}
		< \zeta,
	\end{equation}
	which confirms the validity of Item (II).
	
	
	\subsubsection{Differentiability}
	Let us recall a classical result regarding the geometrical criterion on the differentiability of a Lipschitz submanifold (see \cite[Lemma 4.2]{BB} for instance).
	\begin{Lemma}\label{difli}
		Let \( Z \) be a Lipschitz submanifold of dimension \( n \). If for every \( z \in Z \), the tangent cone \( TC_Z(z) \) is contained in an \( n \)-dimensional space \( L(z) \), then \( Z \) is a differentiable submanifold with \( T_zZ = L(z) \). If moreover \( z \mapsto L(z) \) is continuous, then \( Z \) is of class \( C^1 \).
	\end{Lemma}
	
	In view of  Item (II) in last subsection, we have (see \cite[Proof of Theorem 1.2]{Ne}) for each $x\in\R$, the set
	\[L(x):=\cap_{k\geq 0}DF^{k}(\pi_1^{-1}(g^{-k}(x)))\cdot (\mathcal{C}_\zeta(g^{-k}(x)))\subseteq \mathcal{C}_\zeta(x)\]
	is a 1-dimensional subspace of $T_{\pi_1^{-1}(x)}\R^2$.
	
	
	By (\ref{cooncd}), there holds for each $x\in\R$,
	\[TC_{\Gamma}(x)=\cap_{k\geq 0}DF^{k}(\pi_1^{-1}(g^{-k}(x)))\cdot (TC_{\Gamma}(g^{-k}(x))).\]
	By construction, we have
	\[\cup_{x\in\R}TC_{\Gamma}(x)\subseteq \cup_{x\in\R}\mathcal{C}_\zeta(x),\]
	which yields from the definition of $L(x)$ and that
	\[\cup_{x\in\R}TC_{\Gamma}(x)\subseteq \cup_{x\in\R}L(x).\]
	By Lemma \ref{difli}, the function $\Psi$ is differentiable, with $T_{\pi_1^{-1}(x)}\Gamma=L(x)$ for each $x\in \R$.
	
	
	
	
	
	\subsubsection{Continuous differentiability}
	We prove that $\Psi$ is continuously differentiable. Given a convergent sequence \( x_n \to x \) in $\mathbb{R}$, we need to show that
	\[
	L(x_n) := T_{\pi_1^{-1}(x_n)}\Gamma
	\]
	converges to
	\[
	L(x) := T_{\pi_1^{-1}(x)}\Gamma
	\]
	in the Hausdorff topology. By compactness of the Grassmannian, it suffices to show that \( L(x) \) is the unique accumulation point of the sequence \( L(x_n) \).
	
	Assume, for the sake of contradiction, that there exists an accumulation point \( L'(x) \neq L(x) \) of the sequence \( L(x_n) \). Note that for each \( n \), we have
	\[
	L(x_n) \subseteq \mathcal{C}_\zeta(x_n),
	\]
	where \( \mathcal{C}_\zeta(x) \) is continuous with respect to \( x \in \mathbb{R} \) in the Hausdorff topology. Taking limits, it follows that
	\[
	L'(x) \subseteq \mathcal{C}_\zeta(x).
	\]
	
	For each \( k \geq 0 \), the continuity of \( DF^{-k} \) with respect to \( x \in \mathbb{R} \) implies that
	\[
	DF^{-k}(\pi_1^{-1}(x)) \cdot L'(x)
	\]
	is an accumulation point of the sequence
	\[
	DF^{-k}(\pi_1^{-1}(x_n)) \cdot L(x_n) = T_{\pi_1^{-1}(g^{-k}(x_n))}\Gamma.
	\]
	Again, by continuity of \( \mathcal{C}_\zeta \), we have:
	\[
	DF^{-k}(\pi_1^{-1}(x)) \cdot L'(x) \subseteq \mathcal{C}_\zeta(g^{-k}(x)).
	\]
	By the definition of \( L(x) \), it follows that:
	\[
	L'(x) \subseteq \bigcap_{k \geq 0} DF^k(\pi_1^{-1}(g^{-k}(x))) \cdot \mathcal{C}_\zeta(g^{-k}(x)) = L(x).
	\]
	Since \( L'(x) \) and \( L(x) \) have the same dimension, we must have \( L'(x) = L(x) \), a contradiction. Therefore, \( L(x) \) is the unique accumulation point of \( L(x_n) \), and we conclude that \( x \mapsto L(x) \) is continuous in the Hausdorff topology.
	
	Thus, $\Psi$ is continuously differentiable.
	
	
	\subsection{Normal hyperbolicity of the invariant graph}In general, for a given \( r \geq 1 \), we consider the \( r \)-normal hyperbolicity of the invariant graph under the assumption that a \( C^r \) invariant graph exists. To this end, we first establish the following two lemmas.
	
	
	
	
	
	
	
	
	\begin{Lemma}\label{lem:universal-g-lower}
		Let $\lambda\in(0,1)$, $\tau\in(\lambda,1)$.
		Suppose $F$ admits a $C^1$ invariant graph $\Gamma=\{(x,\Psi(x))\mid x\in\mathbb{T}\}$,
		If the perturbation satisfies
		\[
		\|\phi'\|_\infty < (1-\tau)\Bigl(1-\frac{\lambda}{\tau}\Bigr),
		\]
		then
		\[
		g'(x) > \tau \qquad \text{for all } x\in\mathbb{T}.
		\]
	\end{Lemma}
	
	\begin{proof}
		Recall the derived Herman--Mather formula:
		\begin{equation}\tag{A'}
			g'(x) + \frac{\lambda}{g'(g^{-1}(x))} = 1 + \lambda + \phi'(x), \quad \forall x \in \mathbb{T}.
		\end{equation}
		
		Because $g$ is an orientation-preserving diffeomorphism of $\mathbb{T}$, its lift satisfies $g(x + 1) = g(x) + 1$, hence $\int_0^1 g'(x) dx = 1$ and $g'(x) > 0$.
		
		Set
		$$ B_\lambda = (1 - \tau)\left(1 - \frac{\lambda}{\tau}\right). $$
		The assumption $\|\phi'\|_\infty < B_\lambda$ implies $\phi'(x) > -B_\lambda$ for all $x$. Thus,
		\begin{equation}\label{12}
			g'(x) + \frac{\lambda}{g'(g^{-1}(x))} > \tau + \frac{\lambda}{\tau}, \quad \forall x \in \mathbb{T}.
		\end{equation}
		
		Assume for contradiction that the set $S = \{x \in \mathbb{T} \mid g'(x) \le \tau\} \neq \emptyset$.
		For any $x \in S$, applying $g'(x) \le \tau$ to (\ref{12}) yields:
		$$
		\tau + \frac{\lambda}{g'(g^{-1}(x))} \ge g'(x) + \frac{\lambda}{g'(g^{-1}(x))} > \tau + \frac{\lambda}{\tau},
		$$
		which implies $g'(g^{-1}(x)) < \tau$. Therefore, $g^{-1}(x) \in S$, yielding $g^{-1}(S) \subseteq S$.
		
		By the continuity of $g'$, the Lebesgue measure $\mu(S) > 0$.
		Using the substitution $x = g(y)$, we have:
		$$
		\mu(g^{-1}(S)) = \int_{g^{-1}(S)} 1 dy = \int_S \frac{1}{g'(g^{-1}(x))} dx.
		$$
		Since $g'(g^{-1}(x)) < \tau$ on $S$, it follows that:
		$$
		\mu(g^{-1}(S)) > \frac{1}{\tau} \mu(S).
		$$
		Because $\tau < 1$ and $\mu(S) > 0$, we obtain $\mu(g^{-1}(S)) > \mu(S)$, contradicting $g^{-1}(S) \subseteq S$.
		Thus, $S$ is empty, proving $g'(x) > \tau$ for all $x \in \mathbb{T}$.
	\end{proof}
	
	\begin{Lemma}\label{lem:nh-iff}Fix $r\geq 1$.
		Let $\lambda\in(0,1)$ and assume that
		$F$
		admits a $C^r$ invariant graph $\Gamma=\{(x,\Psi(x))\mid x\in\mathbb{T}\}$.
		Then $\Gamma$ is $r$-normally hyperbolic for $F$ if and only if
		\[
		g'(x)>\lambda^{\frac{1}{r+1}}\qquad\forall x\in\mathbb{T}.
		\]
		where $g(x)=x+\alpha_1+\lambda\Psi(x)+\phi(x)$.
	\end{Lemma}
	
	\begin{proof}
		We prove the two implications separately.
		
		\noindent\textbf{Necessity ($\Rightarrow$).}
		Assume $\Gamma$ is $r$-normally hyperbolic. Differentiating the invariance equation $F(x,\Psi(x))=(g(x),\Psi(g(x)))$ with respect to $x$ gives
\[
DF(x,\Psi(x))\begin{pmatrix} 1 \\ \Psi'(x) \end{pmatrix}
= \begin{pmatrix} g'(x) \\ \Psi'(g(x))g'(x) \end{pmatrix}
= g'(x)\begin{pmatrix} 1 \\ \Psi'(g(x)) \end{pmatrix}.
\]
Thus the tangent vector field $e_1(x)=(1,\Psi'(x))^{\mathsf{T}}$ satisfies
$DF e_1 = g'\, e_1\circ g$, and the tangential multiplier is $\tau(x)=g'(x)>0$. It remains to compute the stable multiplier.

\medskip

For each $x\in\mathbb{T}$ consider the tangent space $T_{(x,\Psi(x))}(\mathbb{T}\times\mathbb{R}) \simeq \mathbb{R}^2$
and its subspace $T_x\Gamma = \operatorname{span}\{(1,\Psi'(x))\}$.  The normal space
$N_x = T_{(x,\Psi(x))}(\mathbb{T}\times\mathbb{R}) / T_x\Gamma$ is one-dimensional.
Define the linear map $\pi_x: \mathbb{R}^2 \to \mathbb{R}$ by
\[
\pi_x(u,v) = v - u\,\Psi'(x).
\]
Clearly $\ker\pi_x = \{(u,v): v=u\Psi'(x)\} = T_x\Gamma$, and $\pi_x$ is surjective.  It follows from the first isomorphism theorem  that $\pi_x$ induces a natural
linear isomorphism $N_x\simeq\mathbb{R}$, still denoted by $\pi_x$,
which sends the equivalence class $[u,v]$ to $v-u\Psi'(x)$.

The derivative $DF(x,\Psi(x))$ maps $T_x\Gamma$ into $T_{g(x)}\Gamma$,
hence it induces a linear map $\overline{DF}(x): N_x \to N_{g(x)}$ characterized by
\[
\overline{DF}(x) \circ \pi_x = \pi_{g(x)} \circ DF(x).
\]
We now compute $\overline{DF}(x)$ in the above trivialization.  Write
\[
DF(x) = \begin{pmatrix}
1+\phi'(x) & \lambda \\
\phi'(x)   & \lambda
\end{pmatrix}.
\]
For any $(u,v)^{\mathsf{T}}$,
\[
DF(x)\begin{pmatrix}u\\ v\end{pmatrix}
= \begin{pmatrix}
(1+\phi'(x))u + \lambda v \\
\phi'(x)u + \lambda v
\end{pmatrix}
=: \begin{pmatrix} U \\ V \end{pmatrix}.
\]
Using the relations obtained from the invariance equation,
\[
g'(x) = 1+\phi'(x)+\lambda\Psi'(x),\qquad
\Psi'(g(x)) = \frac{\phi'(x)+\lambda\Psi'(x)}{g'(x)},
\]
a direct computation yields
\begin{align*}
\pi_{g(x)}(U,V) &= V - U\,\Psi'(g(x)) \\
&= \bigl(\phi'u+\lambda v\bigr) - \bigl((1+\phi')u+\lambda v\bigr)\Psi'(g(x)) \\
&= \Bigl[\phi' - (1+\phi')\Psi'(g(x))\Bigr]u + \lambda\bigl[1-\Psi'(g(x))\bigr]v \\
&= -\frac{\lambda\Psi'(x)}{g'(x)}\,u + \frac{\lambda}{g'(x)}\,v
 = \frac{\lambda}{g'(x)}\bigl(v-u\Psi'(x)\bigr) \\
&= \frac{\lambda}{g'(x)}\,\pi_x(u,v).
\end{align*}
Therefore $\overline{DF}(x)$ acts as multiplication by $\lambda/g'(x)$ on $N_x\simeq\mathbb{R}$.

\medskip

By definition of $r$-normal hyperbolicity, there exists a $DF$-invariant continuous splitting
\[
T_\Gamma(\mathbb{T}\times\mathbb{R}) = T\Gamma \oplus E^s,
\]
where $E^s$ is one-dimensional and $DF$-invariant with uniform contraction.  For each $x$,
the restriction of the projection $\pi_x$ to the stable fiber $E^s_x$,
\[
\pi_x|_{E^s_x} : E^s_x \longrightarrow N_x,
\]
is a linear isomorphism because $E^s_x \cap T_x\Gamma = \{0\}$ and both spaces are
one-dimensional.  Hence we can select a unique continuous section $w(x)\in E^s_x$ such that
\[
\pi_x(w(x)) = 1 \quad \text{for all } x\in\mathbb{T}.
\]
Write $DF(x)w(x) = \nu(x)\, w(g(x))$, where $\nu(x)$ is the stable multiplier.
Applying $\pi_{g(x)}$ to this identity and using the commutativity
$\pi_{g(x)}\circ DF(x) = \overline{DF}(x)\circ\pi_x$, we obtain
\[
\frac{\lambda}{g'(x)} = \overline{DF}(x)(1)
= \overline{DF}(x)\bigl(\pi_x(w(x))\bigr)
= \pi_{g(x)}\bigl(DF(x)w(x)\bigr)
= \nu(x)\,\pi_{g(x)}(w(g(x)))
= \nu(x).
\]
Thus the stable multiplier is precisely $\nu(x)=\lambda/g'(x)$.

\medskip

The definition of $r$-normal hyperbolicity  implies that there exists a
Riemannian metric on $TM|_\Gamma$ such that for every $x$,
\[
\|DF(x)|_{E^s_x}\| \; \bigl\|(DF(x)|_{T_x\Gamma})^{-1}\bigr\|^{\,r} < 1,
\qquad
\|DF(x)|_{E^s_x}\| < 1.
\]
Because both $T_x\Gamma$ and $E^s_x$ are one-dimensional, the operator norm of a linear
map on a one-dimensional space equals the absolute value of the corresponding multiplier.
Consequently,
\[
|\nu(x)| \cdot |\tau(x)|^{-r} < 1,
\]
i.e.\
\[
\frac{\lambda}{g'(x)} \Bigl(\frac{1}{g'(x)}\Bigr)^{\!r} < 1.
\]
This gives $g'(x)^{r+1} > \lambda$.

		
		\medskip
		\noindent\textbf{Sufficiency ($\Leftarrow$).}
		Now assume $g'(x)>\lambda^{1/(r+1)}$ for every $x$.
		Set $\tau=\lambda^{1/(r+1)}$ and $\theta(x)=\lambda/g'(x)$. Then
		$\theta(x)<\lambda/\tau\le \tau< g'(x)$ and $\theta(x)\,g'(x)^{-r}<1$.
		
		\medskip
	
		Take the basis $e_1(x)=(1,\Psi'(x))^{\mathsf{T}}$, $e_2(x)=(0,1)^{\mathsf{T}}$.
		Direct computation gives
		\[
		DF(x)e_1(x)=g'(x)e_1(g(x)),\qquad
		DF(x)e_2(x)=\lambda e_1(g(x))+\frac{\lambda}{g'(x)}e_2(g(x)).
		\]
		Hence the matrix of $DF(x)$ in $(e_1,e_2)$ is
		\[U(x)=\begin{pmatrix} g'(x) & \lambda \\ 0 & \theta(x) \end{pmatrix}.\]
		Look for a stable vector of the form $w(x)=\tilde a(x)e_1(x)+e_2(x)$ with
		$DF(x)w(x)=\theta(x)w(g(x))$. This yields the functional equation
		\[
		\tilde \alpha(g(x))=\frac{(g'(x))^2}{\lambda}\tilde \alpha(x)+g'(x)
		\;\Longleftrightarrow\;
		\tilde \alpha(x)=\frac{\lambda}{(g'(x))^2}\tilde \alpha(g(x))-\frac{\lambda}{g'(x)}.
		\]
		Note $\tau=\lambda^{1/(r+1)}$, $\lambda/\tau^2=\lambda^{1-2/(r+1)}\le1$. Because \[\displaystyle\kappa:=\sup_x\frac{\lambda}{(g'(x))^2}
		<\frac{\lambda}{\tau^2}\le1 \quad \text{for all } r\ge 1, \]
		the operator
		\[(\mathcal{Q}\tilde{\alpha})(x)=\frac{\lambda}{(g'(x))^2}\tilde{\alpha}(g(x))-\frac{\lambda}{g'(x)}\]
		is a strict contraction on $C^0(\mathbb{T})$.
		By the Banach fixed-point theorem there exists a unique continuous $\tilde \alpha$
		solving the equation.
		Define $E^s_x=\operatorname{span}\{w(x)\}$ with $w(x)=\tilde \alpha(x)e_1(x)+e_2(x)$.
		Then $TM|_\Gamma=T\Gamma\oplus E^s$ is a $DF$-invariant splitting and
		$DF|_{E^s_x}$ is multiplication by $\theta(x)$.
		
		\medskip

		The vector fields $e_1(x)$ and $w(x)$ form a continuous frame of the tangent bundle
$TM|_\Gamma$.  We define a continuous
Riemannian metric by requiring $\{e_1(x), w(x)\}$ to be orthonormal.  Concretely, for
any $u,v\in T_{(x,\Psi(x))}M$ written uniquely as
$u = a_1 e_1(x) + a_2 w(x)$,\; $v = b_1 e_1(x) + b_2 w(x)$, set
\[
\langle u,v\rangle_x = a_1 b_1 + a_2 b_2.
\]
This form is symmetric, positive definite, and varies continuously with $x$ because the
frame depends continuously on $x$ and the coordinate expressions are continuous linear
isomorphisms on each fibre.  In this metric,
\begin{align*}
&\|DF|_{T_x\Gamma}\| = \|DF(x) e_1(x)\| = \|g'(x) e_1(g(x))\| = g'(x),\\
&\|DF|_{E^s_x}\| = \|DF(x) w(x)\| = \|\theta(x) w(g(x))\| = \theta(x),\\
&\|(DF|_{T_x\Gamma})^{-1}\| = \bigl\| \tfrac{1}{g'(x)} e_1(x) \bigr\| = \frac{1}{g'(x)}.
\end{align*}
		Thus
		\[
		\|DF|_{E^s_x}\|\;\bigl\|(DF|_{T_x\Gamma})^{-1}\bigr\|^r
		= \frac{\lambda}{(g'(x))^{r+1}} < 1,\qquad
		\|DF|_{E^s_x}\| = \frac{\lambda}{g'(x)} < 1,
		\]
		uniformly in $x$. This means $\Gamma$ is $r$-normally hyperbolic.
	\end{proof}
	\begin{Remark}
  In the proof of the necessity part of Lemma~\ref{lem:nh-iff}, we introduce the quotient bundle $N_x = T_{(x,\Psi(x))}M/T_x\Gamma$ and the induced map $\overline{DF}(x)$, instead of working directly with the normal bundle. The reason for this approach is that the invariant splitting provided by normal hyperbolicity is linked to the computable quotient map through the natural isomorphism $\pi_x$, so the stable multiplier is obtained directly from the commutative relation $\overline{DF}\circ\pi_x = \pi_{g(x)}\circ DF$. Moreover, the quotient trivialisation $\pi_x(u,v)=v-u\Psi'(x)$ is intrinsic to the graph structure without introducing extra choices such as a Riemannian metric.
\end{Remark}
	It suffices to set  \( \tau = \sqrt{\lambda} \) in Lemma~\ref{lem:universal-g-lower} and \( r = 1 \) in Lemma~\ref{lem:nh-iff} to obtain the desired conclusion.
	\begin{Corollary}\label{lexx2}
		Given \(\lambda \in (0,1)\), for any \(\alpha \in \mathbb{R}^2\) and any \(C^1\) perturbation \(\phi\) with zero average satisfying
		\begin{equation}\label{m22conx}
			\|\phi\|_{C^1} < (1-\sqrt{\lambda})^2,
		\end{equation}
		the  \(C^1\) invariant graph persisted by the map \(F\)  is \(1\)-normally hyperbolic.
	\end{Corollary}
	
	
	
	\vspace{1em}
	
	\section{\sc Proof of Theorem \ref{Th1}}
	Let us recall the  Herman--Mather formula
	\begin{equation}\tag{A}
		g(x) + \lambda g^{-1}(x) = (1+\lambda)x + (1-\lambda)\alpha_1 + \lambda\alpha_2 + \phi(x) \quad \forall x \in \mathbb{R},
	\end{equation}
	where
	\[
	g(x) = x + \alpha_1 + \lambda \Psi(x) + \phi(x),
	\]
	and the derived Herman--Mather formula
	\begin{equation}\tag{A'}
		g'(x) + \frac{\lambda }{g'(g^{-1}(x))} = 1+\lambda + \phi_\lambda'(x).
	\end{equation}
	Since \(\phi\) is \(C^1\), in order to ensure that \(\Psi\) has a non-differentiable point, it suffices to guarantee that \(g\) has non-differentiable points. Specifically, it suffices to prove the following proposition.
	
	\begin{Proposition}\label{hf1}
		For any \(\lambda\in(0,1)\) and any rational number \(\beta\), there exist \(\alpha_1,\alpha_2\in\mathbb{R}\) and a non-differentiable \(g\in D^0(\mathbb{T})\) with \(\rho(g)=\beta\) such that, setting
		\[
		G(x) := g(x) + \lambda g^{-1}(x), \quad \forall x \in \mathbb{R},
		\]
		the function \(G\) is of class \(C^\infty\) on \(\mathbb{R}\) and satisfies
		\begin{equation}\label{Gcond}
			\sup_{x\in\mathbb{R}} G'(x) = 1+\lambda + (1-\sqrt{\lambda})^2.
		\end{equation}
	\end{Proposition}
	
	\vspace{1ex}
	\begin{Remark}
		According to \eqref{gdiff}, the derivative of \(\phi\) can be obtained from \(g\). Since \(\int_{0}^{1} \phi(x)\,dx = 0\), the function \(\phi\) is uniquely determined by its derivative. Furthermore, from \eqref{eq:main} we can determine
		\[
		\tilde{\alpha} := (1-\lambda)\alpha_1 + \lambda\alpha_2.
		\]
		Hence \(\alpha_1\) and \(\alpha_2\) are not unique. For simplicity, one can take \(\alpha_2 = 0\), in which case the corresponding dissipative twist map is an exact conformally symplectic twist map.
	\end{Remark}
	
	The remainder of the section is devoted to the proof of Proposition~\ref{hf1}. The proof consists of two parts. In the first part, we construct a \(C^1\) function \(G\) satisfying all the other requirements of Proposition~\ref{hf1}. In the second part, we refine the argument of the first part to construct a \(C^\infty\) function \(G\). To make the construction more quantitative, for any \(s > 1\), we  provide a construction of \(G\) in the Gevrey-\(s\) class.
	
	\subsection{The \(C^1\) case}
	The concrete construction is carried out in the following six steps.
	
	\vspace{1ex}
	
	\textit{Step 1. Constants and notations}
	
	For fixed $\lambda \in (0,1)$, we define
	\begin{equation}
		M_{\lambda}:=1+\lambda+(1-\sqrt{\lambda})^2 = 2+2\lambda-2\sqrt{\lambda}.
	\end{equation}
	
	Consider the quadratic equation $z+\lambda/z = M_{\lambda}$, i.e. $z^2-M_{\lambda}z+\lambda=0$. Its discriminant is $\Delta = M_{\lambda}^2-4\lambda >0$, and the two positive roots are
	\[
	a:=\frac{M_{\lambda}+\sqrt{\Delta}}{2},\quad b:=\frac{M_{\lambda}-\sqrt{\Delta}}{2}.
	\]
	
	One checks directly that $a>1>\sqrt\lambda>b>0$, $ab=\lambda$, and
	\begin{equation}
		a+\frac{\lambda}{a}=b+\frac{\lambda}{b}=M.
	\end{equation}
	
	In the sequel the numbers $a$ and $b$ will serve as the maximal and minimal possible values of the derivative of the homeomorphism $g$ that we are going to build.
	
	\vspace{1ex}
	
	\textit{Step 2. Partition induced by a rational rotation.}
	
	Let $\beta=p/q$ in lowest terms. Assume $q\geq 2$ (the case $q=1$ is trivial and can be obtained analogously). Let $R(x)=x+p/q$ be the rigid rotation on the circle $\mathbb{T}=\mathbb{R}/\mathbb{Z}$ and $\tilde{R}$ be its lift. The orbit of $0$ under $R$ consists of the points
	\[
	\left\{\left\{\frac{kp}{q}\right\} : k=0,1,\dots,q-1 \,\right\}.
	\]
	Then we order them increasingly and denote them by
	\[
	0=y_0<y_1<\dots<y_{q-1}<1.
	\]
	Set $y_q:=y_0+1=1$ for convenience. Define the lengths $L_k:=y_{k+1}-y_k$ for $k=0,\dots,q-1$; clearly $\sum_{k=0}^{q-1}L_k=1$. By construction, for each $k$ there exists a unique $\sigma(k)\in\{0,\dots,q-1\}$ and an integer $m_k\in\{0,1\}$ such that
	\begin{equation}
		\tilde{R}(y_k)=y_k+\frac{p}{q}=y_{\sigma(k)}+m_k.
	\end{equation}
	Note that $\sigma$ is a $q$-cycle. Moreover, the sum of the $m_k$ equals $p$. A crucial observation is that $L_k = L_{\sigma(k)}$ for all $k$; this follows from the fact that $R$ is an isometry, so it maps the interval $[y_k,y_{k+1})$ onto an interval of the same length, which must be some $[y_{\sigma(k)},y_{\sigma(k)+1})$.
	
	\vspace{1ex}
	
	\textit{Step 3. A smooth template function}
	
	\begin{Lemma} \label{lemma1}
		There exists a smooth strictly increasing function $\varphi_c:[0,1]\to[0,1]$ such that
		\[
		\varphi_c(0)=0, \varphi_c(1)=1, and \int_0^1\varphi_c(x)\,dx = c \text{ for any prescribed } c\in(0,1)
		\]
	\end{Lemma}
	\begin{proof}
		Consider the set \[\mathcal{H}=\{\varphi\in C^\infty([0,1]): \varphi(0)=0,\ \varphi(1)=1,\ \varphi'>0\text{ on }(0,1)\}\] endowed with the $C^1$ topology. The map $\Phi:\mathcal{H}\to(0,1)$ defined by $\Phi(\varphi)=\int_0^1\varphi$ is continuous. Take $\varphi_n(x)=x^n$; then $\Phi(\varphi_n)=1/(n+1)\to0$. Take $\varphi_n(x)=1-(1-x)^n$; then $\Phi(\varphi_n)=1-1/(n+1)\to1$. Since $\mathcal{H}$ is connected (it is convex), containing values arbitrarily close to $0$ and $1$, and $\Phi(\phi)$ can't attain $0$ and $1$, hence $\Phi(\mathcal{H})=(0,1)$. Therefore for any $c\in(0,1)$ there exists $\varphi_c\in\mathcal{H}$ with $\Phi(\varphi_c)=c$.
	\end{proof}
	By Lemma \ref{lemma1}, we can construct a smooth strictly decreasing function $\psi:[0,1]\to[b,a]$ such that
	\[
	\psi(0)=a,\quad \psi(1)=b,\quad \int_0^1\psi(t)\,dt=1.
	\]
	
	More precisely, take $\psi(t)=a+(b-a)\varphi_c(t)$ with $0<c<1$ to be determined. Then
	\[
	\int_0^1\psi(t)\,dt = a + (b-a)c.
	\]
	Setting this equal to $1$ gives
	\[
	(b-a)c=1-a\quad\Longrightarrow\quad c=\frac{1-a}{b-a}.
	\]
	Since $a>1>b$, $0<\frac{1-a}{b-a}<1$, so solving for $c$ yields $0<c<1$. Thus $\psi$ is well defined, $C^\infty$, strictly decreasing, and satisfies the required properties.
	
	\vspace{1ex}
	
	\textit{Step 4. Construction of $g$ and its derivative}
	
	On each interval $[y_k,y_{k+1})$ define
	\begin{equation}
		g'(y):=\psi\!\left(\frac{y-y_k}{L_k}\right),\qquad y\in[y_k,y_{k+1}).
	\end{equation}
	Then set
	\begin{equation}
		g(y):=\frac{p}{q}+\int_0^y \tilde{g}(t)\,dt\qquad\text{for }0\le y\le1,
	\end{equation}
	and extend $g$ to $\mathbb{R}$ by $g(y+1)=g(y)+1$. Because $\int_0^1 g'(t)dt=\sum_k L_k\cdot1=1$, we have $g(1)=g(0)+1$. The function $g$ is strictly increasing and continuous. On each open interval $(y_k,y_{k+1})$, the function $g$ is of class $C^\infty$. At the partition points $y_k$ the left derivative equals $b$ and the right derivative equals $a$; since $a\neq b$, $g$ is not differentiable at $y_k$. Hence $g\notin C^1$.
	
	\vspace{1ex}
	
	\textit{Step 5. Rotation number of $g$}
	
	From the definition we have $g(y_k)=y_{\sigma(k)}+m_k$ (this follows by integrating $g'$ over $[0,y_k]$ and using that each subinterval contributes exactly its length $L_i$). Because $\sigma$ is a $q$-cycle and the map sends intervals to intervals of the same length, iterating $q$ times gives
	\[
	g^q(y_k)=y_k+\sum_{i=0}^{q-1}m_{\sigma^i(k)}.
	\]
	The sum $\sum_{i=0}^{q-1}m_{\sigma^i(k)}$ is independent of $k$ and equals $p$. In particular $g^q(0)=p$. Therefore the lift $g$ has rotation number $\rho(g)=p/q$.
	
	\vspace{1ex}
	
	\textit{Step 6. Differentiability of $G$ and the sharp upper bound for $G'$}
	
	Because $g$ is strictly increasing and bijective, $g$ has an inverse, denoted by $g^{-1}$. Then $g^{-1}$ is also a homeomorphism and satisfies $g^{-1}(x+1)=g^{-1}(x)+1$. Define
	\[
	G(x):=g(x)+\lambda g^{-1}(x),\quad \text{for } x\in \mathbb{R}.
	\]
	
	Let $x\in\mathbb{R}$. There exists $y\in[0,1)$ and $n\in\mathbb{Z}$ such that $x=g(y)+n$. Since $g'$ is $1$-periodic and $(g^{-1})'(x)=1/g'(g^{-1}(x))$, we obtain
	\[
	G'(x)=g'(x)+\lambda(g^{-1})'(x)=g'(g(y))+\frac{\lambda}{g'(y)}.
	\]
	
	Now we prove a useful inequality.
	
	\begin{Lemma}\label{lemma2}
		For the template function $\psi$ constructed in Step~3, one has
		\begin{equation}
			\int_0^u\psi(t)\,dt\ge u\qquad\text{for all }u\in[0,1].
		\end{equation}
	\end{Lemma}
	
	\begin{proof}
		Define $\Phi_*(u)=\int_0^u\psi(t)\,dt-u$. Then $\Phi_*(0)=0$ and $\Phi_*'(u)=\psi(u)-1$. Since $\psi$ is strictly decreasing with $\psi(0)=a>1$ and $\psi(1)=b<1$, there exists a unique $c\in(0,1)$ such that $\psi(c)=1$. Hence $\Psi$ increases on $[0,c]$ and decreases on $[c,1]$. Because $\Phi_*(1)=\int_0^1\psi-1=0$ and $\Phi_*(0)=0$, we conclude $\Phi_*(u)\ge0$ for all $u$.
	\end{proof}
	
	\vspace{1ex}
	
	Now take any $y$ in some interval $(y_k,y_{k+1})$ and set $u=(y-y_k)/L_k\in(0,1)$. Then $g'(y)=\psi(u)$. Using the definition of $g$,
	\[
	g(y)=m_k+y_{\sigma(k)}+\int_{y_k}^y\psi\!\left(\frac{t-y_k}{L_k}\right)dt
	=m_k+y_{\sigma(k)}+L_k\int_0^u\psi(t)\,dt.
	\]
	Because $L_k=L_{\sigma(k)}$, the point $g(y)$ lies in the interval $[m_k+y_{\sigma(k)},m_k+y_{\sigma(k)+1})$ and its normalised coordinate is $\int_0^u\psi(t)dt$. Consequently,
	\[
	g'(g(y))=\psi\!\left(\int_0^u\psi(t)dt\right).
	\]
	
	By the Lemma \ref{lemma2}, $\int_0^u\psi\ge u$, and since $\psi$ is decreasing we obtain
	\[
	g'(g(y))=\psi\!\left(\int_0^u\psi\right)\le\psi(u)=g'(y).
	\]
	This inequality also holds at the partition points by continuity of the left and right hand sides. Therefore for every $y\in[0,1]$ we have $g'(g(y))\le g'(y)$.
	
	Using this,
	\[
	G'(x)\le g'(y)+\frac{\lambda}{g'(y)}.
	\]
	Since $g'(y)\in[b,a]$, and the function $h(t)=t+\lambda/t$ attains its maximum on $[b,a]$ at the endpoints with $h(a)=h(b)=M_{\lambda}$, we conclude
	\[
	G'(x)\le M_{\lambda}\qquad\text{for all }x\in\mathbb{R}.
	\]
	
	Finally, consider $x=y_0=0$. Then $g(0)=p/q$, $g'(0^+)=a$, and $g'(g(0))=g'(\frac{p}{q})=a$, so
	\[
	G'(0^+)=a+\frac{\lambda}{a}=M_{\lambda}.
	\]
	The left-hand limit gives the same value. Hence $G'$ is continuous at $0$ and actually everywhere, because at every partition point the left and right limits of $G'$ equal $M_{\lambda}$. Thus $G\in C^1(\mathbb{R})$ and $\sup G' = M_{\lambda}$.
	
	This completes the proof of the $C^1$ case.
	
	
	\subsection{From \(C^1\) to Gevrey regularity}
	
	The construction in Proposition~\ref{hf1} used a smooth template function $\psi$ with $\psi(0)=a$, $\psi(1)=b$ and $\int_0^1\psi=1$, but no further conditions at the endpoints. Consequently the resulting $G$ was only proved to be $C^1$. To obtain Gevrey regularity for $G$ for any exponent $s>1$ we replace the function space for the auxiliary function $\varphi$ (which defines $\psi=a+(b-a)\varphi$) by one whose elements are flat at the endpoints in the Gevrey sense.
	
	Recall that a \(C^\infty\) function \(f\) on an interval belongs to the Gevrey class \(G^s\) (\(s>1\)) if for every compact set \(K\) there exist constants \(C,R>0\) such that
	\[
	\sup_{x\in K}|f^{(k)}(x)|\le C\, k!^{\,s}R^{-k}\qquad\forall k\ge0.
	\]
	
	For \(s\leq 1\) this gives real analyticity; for \(s>1\) the class is strictly larger than analytic and still closed under composition, inversion, algebraic operations when the involved functions are bounded away from zero (see for instance \cite[Appendix B]{BF}).
	
	%\vspace{1ex}
	%	
	%	\begin{Corollary}
		%		For any \(\lambda\in(0,1)\) and any rational number \(\beta\) and any fixed $s>1$, there exist \(\alpha_1,\alpha_2\in\mathbb{R}\) and a non-differentiable \(g\in D^0(\mathbb{T})\) with \(\rho(g)=\beta\) such that, setting
		%		\[
		%		G(x) := g(x) + \lambda g^{-1}(x), \quad \forall x \in \mathbb{R},
		%		\]
		%		the function \(G\) is of class \(G^s\) on \(\mathbb{R}\) and satisfies
		%		\begin{equation}\label{Gcond}
			%			\sup_{x\in\mathbb{R}} G'(x) = 1+\lambda + (1-\sqrt{\lambda})^2.
			%		\end{equation}
		%		
		%		Moreover, such a construction is impossible for $s=1$ (i.e. $G$ cannot be real analytic).
		%	\end{Corollary}
	
	\vspace{1ex}
	
	
	We modify Step~3 of the proof of the \(C^1\) case.
	
	\vspace{1ex}
	
	\subsubsection{Gevrey flat template}
	\begin{Lemma}
		Let \(s>1\). Define
		\[
		\xi_s(x):=
		\begin{cases}
			\exp\!\left(-\bigl(x(1-x)\bigr)^{-\frac{1}{s-1}}\right), & x\in(0,1),\\
			0, & x=0,1.
		\end{cases}
		\]
		Then the following hold:
		\begin{enumerate}
			\item[(i)] \(\xi_s\in C^\infty([0,1])\), \(\xi_s(x)>0\) for \(x\in(0,1)\), and
			\[
			\xi_s^{(k)}(0)=\xi_s^{(k)}(1)=0\qquad\forall k\ge0.
			\]
			\item[(ii)] \(\xi_s\in G^s([0,1])\), i.e. there exist constants \(C,R>0\) such that
			\[
			\sup_{x\in[0,1]}|\xi_s^{(k)}(x)|\le C\, k!^{\,s}R^{-k}\qquad\forall k\ge0.
			\]
			
		\end{enumerate}
	\end{Lemma}
	
	\begin{proof}
		We prove each part in turn.
		
		\medskip
		\noindent\textit{Proof of (i).} Note that $\xi_s$ is positive and analytic on $(0,1)$, which yields that $\xi_s\in C^\infty((0,1))$. By symmetry it suffices to examine the behaviour near \(x=0\).
		
		Define \(H(x)=e^{-x^{-\frac{1}{s-1}}}\) for \(x>0\), then \[\xi_s(x)=H(x)H(1-x).\]Set $t=x^{-\frac{1}{s-1}}$, then $x=t^{-(s-1)}$ and \[x^{-m}H(x)=t^{(s-1)m}e^{-t}\text{ for any integer } m\ge0.\]  Since $e^t$ increases faster than any power of $t$ as $t$ goes to infinity, so we have \[\lim_{t\to \infty}t^{(s-1)m}e^{-t}=0,\]which yields that \[\lim_{x\to0^+} x^{-m}H(x)=0.\]
		
		By induction, \(H^{(k)}(x)\) is a finite linear combination of terms of the form \(x^{-n}H(x)\) with \(n\ge0\). Hence \[H^{(k)}(0^+)=0\text{ for every } k.\]Since $H(x)$ is analytic at $x=1$, by the Cauchy estimate, we obtain that \[H^{(k)}(1)=0\text{ for every } k.\]
		
		Therefore, by the Leibniz rule and symmetry, we obtain that \[\xi_s^{(k)}(0)=0,\; \xi_s^{(k)}(1)=0\text{ for every } k.\]
		
		\medskip
		\noindent\textit{Proof of (ii).} We only need to derive the Gevrey estimate near \(0\); near \(1\) it follows by symmetry.
		
		The function \(H(z)=e^{-z^{-\frac{1}{s-1}}}\) is holomorphic in the sector \(\Sigma_\theta=\{z\in\mathbb C: |\arg z|<\theta\}\) for any \(\theta<\pi (s-1)/2\). Fix such a \(\theta\) and choose a constant \(c_\theta>0\) such that for \(z\) with \(|z-x|\le x/2\) and \(x\in(0,1/2)\), we have $\Re(z^{-\frac{1}{s-1}})\ge c_\theta x^{-\frac{1}{s-1}}$, then \[|H(z)|\leq e^{-c_\theta x^{-\frac{1}{s-1}}}.\]
		For \(x\in(0,1/2)\), let \(\gamma_x\) be the circle centered at \(x\) with radius \(r=x/2\). By Cauchy's integral formula,
		\[
		|H^{(k)}(x)| = \left|\frac{k!}{2\pi i}\int_{\gamma_x}\frac{H(z)}{(z-x)^{k+1}}\,dz\right|
		\le \frac{k!}{r^k}\sup_{z\in \gamma_x}|H(z)|
		\le k!\left(\frac{2}{x}\right)^k e^{-c_\theta x^{-\frac{1}{s-1}}}.
		\]
		Set \(t=c_\theta x^{-\frac{1}{s-1}}\). Then \(x^{-k}=(t/c_\theta)^{k(s-1)}\) and
		\[
		|H^{(k)}(x)| \le 2^k c_\theta^{-k(s-1)} k! \, t^{k(s-1)} e^{-t}.
		\]
		The supremum of \(t^{k(s-1)}e^{-t}\) over \(t>0\) is \((k(s-1))^{k(s-1)}e^{-k(s-1)}\le C_1^k k!^{s-1}\) for some \(C_1>0\). Therefore
		\[
		|H^{(k)}(x)| \le C_2^k k!\, k!^{\,s-1} = C_2^k k!^{\,s},
		\]
		with \(C_2>0\) depending on \(\theta,s\). This proves the Gevrey estimate on \((0,1/2)\). The same estimate on \((1/2,1)\) gives the global bound on \([0,1]\) (with possibly different constants, which can be combined by taking the maximum). Hence \(H\in G^s([0,1])\). Since $\xi_s(x)=H(x)H(1-x)$, so $\xi_s \in G^s([0,1])$.
	\end{proof}
	
	Normalising the integral of $\xi_s$ gives
	\[
	\varphi_0(x)=\frac{\int_0^x\xi_s(t)\,dt}{\int_0^1\xi_s(t)\,dt}.
	\]
	Then \(\varphi_0(0)=0,\ \varphi_0(1)=1,\ \varphi_0'(x)=\frac{\xi_s(x)}{\int_0^1\xi_s(t)\,dt}\), and since \(\xi_s\) is flat at the endpoints, \[\varphi_0^{(k)}(0)=\frac{\xi_s^{(k-1)}(0)}{\int_0^1\xi_s(t)\,dt}=0\] and similarly at \(1\) for all \(k\ge1\). The Gevrey regularity of \(\varphi_0\) follows from that of \(\xi_s\). Thus \(\varphi_0\in G^s([0,1])\) and satisfies the endpoint derivative conditions. Let \(c_0:=\int_0^1\varphi_0\).
	
	To obtain an element with an arbitrarily prescribed integral \(c\in(0,1)\), one use convex combinations and a rescaling argument.
	Let \(\varphi_1(x)=1-\varphi_0(1-x)\); then \(\varphi_1\in G^s([0,1])\) shares the same flatness properties and satisfies \(\int_0^1\varphi_1=1-c_0\).
	Convex combinations of \(\varphi_0\) and \(\varphi_1\) realise every integral value between \(c_0\) and \(1-c_0\).
	For $\varepsilon\in(0,1)$, set
	\[
	\varphi_{0,\varepsilon}(x):=\begin{cases}
		\varphi_0(\frac{x}{\varepsilon}), & x\in(0,\varepsilon),\\
		0, & x\in(\varepsilon,1).
	\end{cases}
	\]
	and set $\varphi_{1,\varepsilon}(x):=1-\varphi_{0,\varepsilon}(1-x)$. We obtain that $\int_{0}^{1}\varphi_{0,\varepsilon}=\varepsilon c_0$, $\int_{0}^{1}\varphi_{1,\varepsilon}=1-\varepsilon c_0$.
	Since $\varepsilon$ can be chose arbitrarily small, the range of attainable integrals can be extended to the whole interval \((0,1)\).
	Consequently, we may fix a function \(\varphi\in G^s([0,1])\) such that
	\[
	\int_0^1\varphi=\frac{1-a}{b-a},
	\]
	and \(\varphi^{(k)}(0)=\varphi^{(k)}(1)=0\) for all \(k\ge 1\).
	Define
	\[
	\psi(t)=a+(b-a)\varphi(t),\qquad t\in[0,1].
	\]
	Then \(\psi\) is strictly decreasing, \(\psi(0)=a\), \(\psi(1)=b\), \(\int_0^1\psi=1\), and \(\psi\in G^s([0,1])\) with \(\psi^{(k)}(0)=\psi^{(k)}(1)=0\) for all \(k\ge 1\).
	
	\vspace{1ex}
	
	\subsubsection{Construction of \(g\) and \(G\)}
	Replace the function \(\psi\) in Steps~2,~4 and~5 of the \(C^1\) case by the new \(\psi\) constructed above.
	The rest of the procedure (the partition \(\{y_k\}\), the definition \(g'(y)=\psi((y-y_k)/L_k)\), etc.) remains unchanged.
	Exactly as before we obtain a circle map \(g\notin C^1\) with rotation number \(p/q\) and \[G(x)=g(x)+\lambda g^{-1}(x)\in C^1.\]
	
	\vspace{1ex}
	
	\subsubsection{Gevrey regularity of \(G\)}
	
	
	On each open interval \((y_k,y_{k+1})\), \(g'\) is a composition of \(\psi\) with an affine map, hence \(g\in G^s\) there. It remains to check the behaviour at the break points \(y_k\).
	
	For \(t>0\) small, define
	\[
	\eta(t):=g'(y_k+t)-a.
	\]
	Since \(g'(y_k+t)=\psi(t/L_k)\) and \(\psi\) is flat at \(0\) in the sense that all derivatives vanish, we have
	\[
	\eta(0)=0,\qquad \eta^{(j)}(0)=0\quad\forall j\ge0,
	\]
	and \(\eta\) satisfies the Gevrey estimates.
	Integrating,
	\[
	g(y_k+t)=g(y_k)+\int_0^t g'(y_k+u)\,du
	= y_{\sigma(k)}+m_k+a t+E(t),
	\]
	where \(E(t):=\int_0^t\eta(u)\,du\). Clearly \(E\) is also Gevrey-flat.
	
	Now, since \(g\) maps break points to break points, we have
	\[
	g'(y_{\sigma(k)}+m_k+z)=a+\eta_0(z),
	\]
	where \(\eta_0\) is the analogous flat remainder on the right of \(y_{\sigma(k)}+m_k\), with the same Gevrey properties. Taking \(z=a t+E(t)\), we get
	\[
	g'(g(y_k+t)) = g'(y_{\sigma(k)}+m_k+a t+E(t))
	= a+\eta_0(a t+E(t)).
	\]
	Because \(a t+E(t)\) is a \(G^s\) function and vanishes at \(t=0\), the composition \(\eta_0(a t+E(t))\) is again Gevrey-flat: its derivatives at 0 are zero and the Gevrey estimates are preserved (composition of Gevrey functions is Gevrey). Hence
	\[
	g'(g(y_k+t))=a+\eta_1(t)
	\]
	with \(\eta_1:=\eta_0(a t+E(t))\in G^s\), \(\eta_1^{(j)}(0)=0\) for all \(j\).
	
	Next, consider the term \(\lambda/g'(y_k+t)=\lambda/(a+\eta(t))\). Since \(\eta(0)=0\) and \(a>0\), for small \(t\) we have \(|\eta(t)|<a/2\). Expand:
	\[
	\frac{\lambda}{a+\eta(t)}
	= \frac{\lambda}{a}\cdot\frac{1}{1+\eta(t)/a}
	= \frac{\lambda}{a}\sum_{j=0}^{\infty}\left(-\frac{\eta(t)}{a}\right)^j.
	\]
	The infinite series converges uniformly for small \(t\). For each \(j\), the function \(\eta(t)^j\) is Gevrey-flat, since the product of flat functions remains flat and Gevrey. The sum is also Gevrey-flat because the Gevrey class is closed under uniform limits of functions with uniform Gevrey bounds. Thus we can write
	\[
	\frac{\lambda}{a+\eta(t)} = \frac{\lambda}{a} + \eta_2(t),
	\]
	where \(\eta_2:=\frac{\lambda}{a} \sum_{j=1}^{\infty}\left(-\frac{\eta(t)}{a}\right)^j\in G^s\) and \(\eta_2^{(j)}(0)=0\) for all \(j\).
	
	Combining the two parts, for \(x=g(y_k+t)+n\), we obtain
	\[
	G'(x)=g'(g(y))+\frac{\lambda}{g'(y)}
	= \bigl(a+\eta_1(t)\bigr)+\left(\frac{\lambda}{a}+\eta_2(t)\right)
	= M_\lambda + \eta_*(t),
	\]
	where \(M_\lambda=a+\lambda/a\) and \(\eta_*:=\eta_1+\eta_2\) is Gevrey-flat.
	
	The same computation from the left side (\(y=y_k-t\)) yields
	\[
	G'(x)=M_\lambda+\tilde\eta_*(t),
	\]
	with \(\tilde\eta_*\) also Gevrey-flat. Hence the left and right derivatives of \(G'\) of all orders at \(x_k\) are equal: indeed, for any \(j\ge1\),
	\[
	(G')^{(j)}(x_k^+)=\eta_*^{(j)}(0)=0,
	\]
	and similarly from the left. Thus \(G'\) is \(C^\infty\) at \(x_k\), and moreover the Gevrey estimates for the derivatives of \(G'\) across the break point are satisfied because the flat remainder \(\eta_*\) belongs to \(G^s\). On each open interval, \(G'\) is \(G^s\) as well. Therefore \(G\in G^s(\mathbb{R})\).
	
	\vspace{1ex}
	
	This completes the proof of Proposition~\ref{hf1}.
	
	
	
	%	\begin{Remark}
		%		For a \(G^s\) perturbation (\(s \leq 1\)), the construction presented here is impossible: an analytic function that is flat at a point must be constant, which would force \(G\) to be constant, contradicting \(\sup G' = M > 1+\lambda\). On the other hand, according to \cite[Corollary 4]{LLW26}, if \(\alpha_1\) is a Diophantine number and \(\alpha_2=0\), then no trigonometric polynomial (of finite degree) can meet the requirements of Theorem~\ref{Th1}. Note that a trigonometric polynomial (of finite degree) corresponds to the \(G^0\) class. Therefore, Theorem \ref{Th1} is optimal in the sense that it excludes the \(G^s\) classes for all \(s \in (0,1]\).
		%	\end{Remark}
	
	
	
	\vspace{1em}
	
	\section{\sc Verification of Remark \ref{R2}}\label{R3}
	Corollary~\ref{lexx2} yields \(A \subseteq B\). The NHIM theorem asserts that \(B \subseteq C\). In general, the converse inclusions are false.
	
	\subsection{\(B \neq A\)}
	
	First, we construct an example showing that \(B \neq A\).
	\begin{Proposition}\label{lem:counterexample}
		For every $\lambda\in(0,1)$ there exists $\alpha\in\mathbb{R}^2$ and
		$\phi_*\in C^1(\mathbb{T})$ with zero average such that
		\begin{itemize}
			\item the map $F_*$ admits a $C^1$ invariant graph
			$\Gamma_*=\{(x,\Psi_*(x))\}$ which is $1$-normally hyperbolic,
			\item $\|\phi_*\|_{C^1} > (1-\sqrt{\lambda})^2$.
		\end{itemize}
	\end{Proposition}
	
	\begin{proof}
		Define $f_*(x):=x+\frac{\lambda}{1-x}-\lambda.$
		
		Since $f_*(x)$ is continuous on $[0,1-\sqrt{\lambda}]$, and
		\[
		f_*(1-\sqrt{\lambda})=1-\lambda>(1-\sqrt{\lambda})^2,\, f_*(0)=0.
		\]
		There exists a $t\in(0,1-\sqrt{\lambda})$ such that
		\[
		f_*(t)=t+\frac{\lambda}{1-t}-\lambda > (1-\sqrt{\lambda})^2.
		\]
		Set
		\[
		a' = 1+t,\qquad b' = 1-t,
		\]
		so that \(a'>1>b'>\sqrt{\lambda}\) and \(a'+b'=2\).
		
		Define a 1-periodic continuous function $p_*:\mathbb{R}\to (0,\infty)$ on $[0,1]$ by:
		\[
		p_*(x):=\begin{cases}
			a', & 0\leq x\leq \dfrac{1}{8}\\[4pt]
			1-8(a'-1)(x-\dfrac{1}{4}), & \dfrac{1}{8}\le x\le \dfrac{1}{4},\\[4pt]
			1, & \dfrac{1}{4}\le x\le \dfrac{3}{8},\\[4pt]
			1-8(1-b')(x-\dfrac{3}{8}), & \dfrac{3}{8}\le x\le \dfrac{1}{2},\\[4pt]
			b', & \dfrac{1}{2}\le x\le \dfrac{5}{8},\\[4pt]
			1+8(1-b')(x-\dfrac{3}{4}), & \dfrac{1}{8}\le x\le \dfrac{1}{4},\\[4pt]
			1, & \dfrac{3}{4}\le x\le \dfrac{7}{8},\\[4pt]
			1+8(a'-1)(x-\dfrac{7}{8}), & \dfrac{7}{8}\le x\le 1.
		\end{cases}
		\]
		Then we obtain \(\int_0^1 p_*(x)\,dx = 1\) and \(p_*(x)\ge b'\) for all $x$.
		
		Define
		\[
		g_*(x) = d + \int_{0}^{x} p_*(t)\,dt,\qquad x\in\mathbb{R},
		\]
		where the constant \(d\) is chosen so that \(g_*(1/2)=0\).  Since
		\(\int_0^1 p_* = 1\), we have \(g_*(x+1)=g_*(x)+1\); thus \(g_*\) is a \(C^1\)
		diffeomorphism of the line covering a circle diffeomorphism.  Its
		derivative satisfies \(g_*'(x)=p_*(x)\ge b'>\sqrt{\lambda}>0\) for all \(x\), so $g_*$ belongs to $C^1$ and has a $C^1$ inverse $g_*^{-1}$.
		
		
		By construction, \(g_*(1/2)=0\), then the preimage of $\zeta_0=0$ under \(g_*\) is
		\(g_*^{-1}(\zeta_0)=1/2\).  Therefore
		\[
		g_*'(\zeta_0)=p_*(0)=a',\qquad
		(g_*^{-1})'(\zeta_0) = \frac{1}{g_*'(g_*^{-1}(\zeta_0))} = \frac{1}{g_*'(1/2)} = \frac{1}{b'}.
		\]
		
		
		Set
		\[
		\widetilde\phi(x) = g_*(x) + \lambda g_*^{-1}(x) - (1+\lambda)x.
		\]
		This is a \(1\)-periodic function.  Let \(\bar\phi = \int_0^1 \widetilde\phi\)
		and define \(\phi_* = \widetilde\phi - \bar\phi\); then \(\int_0^1 \phi_*=0\) and
		\(\phi_*'=\widetilde\phi'\).
		At the point \(\zeta_0=0\) we have
		\[
		\phi_*'(0) = a' + \frac{\lambda}{b'} - (1+\lambda) = t + \frac{\lambda}{1-t} - \lambda = f_*(t).
		\]
		By the choice of \(t\), \(f_*(t) > (1-\sqrt{\lambda})^2\), so
		\[\|\phi_*\|_{C^1}= \|\phi_*'\|_\infty > (1-\sqrt{\lambda})^2.\]
		
		
		Choose \(\alpha_1\) arbitrarily (e.g. \(\alpha_1=0\)) and pick \(\alpha_2\) so that
		\(\bar\phi = (1-\lambda)\alpha_1 + \lambda\alpha_2\).  Then the Herman-Mather formula
		\[
		g_*(x) + \lambda g_*^{-1}(x) = (1+\lambda)x + (1-\lambda)\alpha_1 + \lambda\alpha_2 + \phi_*(x)
		\]
		holds.  Defining \(\Psi_*(x) = \frac{1}{\lambda}\bigl(g_*(x)-x-\alpha_1-\phi_*(x)\bigr)\)
		yields a \(C^1\) invariant graph \(\Gamma_*=\{(x,\Psi_*(x))\}\).
		
		
		Because \(g_*'(x)\ge b'>\sqrt{\lambda}\), by Corollary \ref{lexx2},  \(\Gamma_*\) is \(1\)-normally
		hyperbolic.
	\end{proof}
	\subsection{\(C \neq B\)}
	
	For general systems, it is straightforward to construct counterexamples showing that \(C \neq B\). However, for the standard-like conformal symplectic twist map considered here, such a counterexample is not trivial to obtain, especially if we require that \(\|\phi\|_{C^1} \to 0\) as \(\lambda \to 1^-\). The following lemma provides an explicit construction.
	
	\begin{Proposition}\label{lexx33}
		For every $\lambda\in(0,1)$, one can choose \(\alpha \in \mathbb{R}^2\) and a \(C^1\) perturbation \(\phi^*\) with zero average satisfying
		\begin{equation}\label{m22conx5}
			\|\phi^*\|_{C^1} \leq 3\sqrt{1-\lambda},
		\end{equation}
		such that the map \(F^*\) admits a unique \(C^1\) invariant graph, but this invariant graph is not \(1\)-normally hyperbolic.
	\end{Proposition}
	
	
	
	
	
	
	
	
	We prove Proposition~\ref{lexx33} under the assumption that the following Lemma holds.
	
	\begin{Lemma}\label{16p}
		For every $\lambda\in(0,1)$ set $\delta=\sqrt{1-\lambda}$.
		There exist constants $\alpha_1,\alpha_2\in\mathbb{R}$ and $C^1$ functions
		$\phi^*,\Psi^*:\T\to\mathbb{R}$ such that
		\begin{itemize}
			\item[(i)] $\|\phi^*\|_{C^1}\le 3\delta$,
			\item[(ii)] $\displaystyle\int_{\T}\phi^*(x)\,dx=0$,
			\item[(iii)] there exists $\xi_0\in\T$ with $\left(\phi^*\right)'(\xi_0)=-\delta$,
			\item[(iv)] for the associated invariant graph $\Psi^*$ by $F^*$ we have
			\[
			|\left(\Psi^*\right)'(\xi_0)|\le\frac{\delta}{1+\delta}.
			\]
		\end{itemize}
	\end{Lemma}
	Recall that if the map admits a $C^1$ invariant graph, then the circle diffeomorphism
	\begin{equation}\label{dfg}
		g^*(x)=x+\alpha_1+\lambda\Psi^*(x)+\phi^*(x)
	\end{equation}
	satisfies the Herman-Mather formula (\ref{eq:main})
	\begin{equation}
		g^*(x)+\lambda (g^*)^{-1}(x)=(1+\lambda)x+(1-\lambda)\alpha_1+\lambda\alpha_2+\phi^*(x).
	\end{equation}
	%It is easy to verify that $\tilde g$ we constructed in Lemma~\ref{16p\} (related to $\phi_\lambda$ and $\Psi$ by (\ref{dfg})) has rotation number $1/2$.
		
		\subsection*{Proof of Proposition~\ref{lexx33}}
		By Corollary \ref{lexx2}, it suffices to show that there exists \(\xi_0\) such that
		\[
		(g^*)'(\xi_0) \leq \sqrt{\lambda}.
		\]
		From the definition of \(g^*\) and the construction in Lemma~\ref{16p}, a direct computation yields
		\[
		(g^*)'(\xi_0) = 1 + (\phi^*)'(\xi_0) + \lambda (\Psi^*)'(\xi_0) \leq 1 - \delta + \frac{\lambda \delta}{1+\delta} = \lambda < \sqrt{\lambda}.
		\]
		Hence, Proposition~\ref{lexx33} is proved.
		\subsection*{Proof of Lemma \ref{16p}}
		
		
		
		Fix $\lambda\in(0,1)$ and let $\delta=\sqrt{1-\lambda}\in(0,1)$.
		Choose a small parameter $\varepsilon=\frac18$ and define
		\[
		a''=\varepsilon\delta,\qquad b''=\frac12-\varepsilon,\qquad c''=\varepsilon\delta^2.
		\]
		
		\vspace{1ex}
		
		\noindent\textit{Construction of a diffeomorphism $g^*$.}
		Define a $1$-periodic continuous function $(g^*)':\mathbb{R}\to\mathbb{R}$ on $[0,1]$ by
		\[
		(g^*)'(x)=
		\begin{cases}
			1+\delta-\dfrac{\delta}{a''}x, & 0\le x\le a'',\\[4pt]
			1, & a''\le x\le b'',\\[4pt]
			1-\dfrac{1-\lambda}{\frac12-b''}(x-b''), & b''\le x\le\frac12,\\[4pt]
			\lambda+\dfrac{1-\lambda}{a''}\Bigl(x-\frac12\Bigr), & \frac12\le x\le\frac12+a'',\\[4pt]
			1, & \frac12+a''\le x\le 1-c'',\\[4pt]
			1+\dfrac{\delta}{c''}\Bigl(x-1+c''\Bigr), & 1-c''\le x\le 1.
		\end{cases}
		\]
		One checks that $(g^*)'$ is strictly positive, continuous, and satisfies
		\[
		\int_0^{1/2} (g^*)'(x)\,dx = \frac12,\qquad
		\int_0^1 (g^*)'(x)\,dx = 1.
		\]
		Set
		\[
		g^*(x)=\frac12+\int_0^x (g^*)'(t)\,dt,\qquad x\in\mathbb{R}.
		\]
		Then $g^*(0)=\frac12$, $ g^*(\frac12)=1$, and $g^*(x+1)=g^*(x)+1$.
		Thus $g^*$ is a $C^1$ orientation-preserving diffeomorphism of $\mathbb{R}$ with
		rotation number $1/2$. In particular $(g^*)^{-1}(x)$ exists and is $C^1$.
		
		\vspace{1ex}
		
		\noindent\textit{Definitions of $\phi^*$ and $\Psi^*$ .}
		Choose $\alpha_1=\frac12$ and define
		\[
		\widetilde\phi(x)=g^*(x)+\lambda (g^*)^{-1}(x)-(1+\lambda)x-\frac{1-\lambda}{2}.
		\]
		Let $c'=\int_0^1\widetilde\phi(x)\,dx$ and set
		\[
		\phi^*(x)=\widetilde\phi(x)-c',\qquad
		\alpha_2=\frac{c'}{\lambda}.
		\]
		Then $\int_{\T}\phi^*=0$.  Finally define
		\[
		\Psi^*(x)=\frac{g^*(x)-x-\alpha_1-\phi^*(x)}{\lambda}.
		\]
		A direct substitution shows that (A) holds.
		
		\vspace{1ex}
		
		\noindent\textit{Verification of (iii) and (iv).}
		Take $\xi_0=\frac12$.  Because $g^*(0)=\frac12$ and $g^*(\frac12)=1$, we have
		$(g^*)^{-1}(\xi_0)=0$.  From the definitions of $(g^*)'$,
		\[
		(g^*)'(\xi_0)=(g^*)'\Bigl(\frac12\Bigr)=\lambda,\qquad
		(g^*)'(0)=1+\delta.
		\]
		Using \[((g^*)^{-1})'(\xi_0)=\frac{1}{(g^*)'((g^*)^{-1}(\xi_0))}=\frac{1}{(g^*)'(0)}=\frac{1}{1+\delta},\] we compute
		\begin{align*}
			\left(\phi^*\right)'(\xi_0) &= (g^*)'(\xi_0)+\lambda((g^*)^{-1})'(\xi_0)-(1+\lambda) \\
			&= \lambda+\frac{\lambda}{1+\delta}-1-\lambda \\
			&= -\delta.
		\end{align*}
		
		
		We have $\left(\phi^*\right)'(\xi_0)=-\delta$ by construction. The definition of $\Psi^*$ gives
		\[
		\left(\Psi^*\right)'(x) = \frac{1}{\lambda}\bigl((g^*)'(x)-1-(\phi^*)'(x)\bigr).
		\]
		Plugging $x=\xi_0$:
		\[
		\left(\Psi^*\right)'(\xi_0) = \frac{1}{\lambda}\bigl((g^*)'(\xi_0)-1-(-\delta)\bigr)
		= \frac{1}{\lambda}(\lambda-1+\delta).
		\]
		Since $\lambda=1-\delta^2$, we have
		\[
		\left(\Psi^*\right)'(\xi_0) = \frac{-\delta^2+\delta}{1-\delta^2}=\frac{\delta}{1+\delta}.
		\]
		Thus condition (iv) holds with equality.
		
		\vspace{1ex}
		
		\noindent\textit{Estimate of the $C^1$-norm.}
		From the construction of $(g^*)'$ we have the uniform bounds
		\[
		|(g^*)'(x)-1|\le \delta,\qquad \min_{x\in\mathbb{T}} (g^*)'(x)=\lambda>0.
		\]
		For the inverse function,
		\[
		|((g^*)^{-1})'(x)-1|
		= \Bigl|\frac{1}{(g^*)'((g^*)^{-1}(x))}-1\Bigr|
		\le \max_{y:=(g^*)^{-1}(x)}\frac{|(g^*)'(y)-1|}{(g^*)'(y)}
		\le \frac{\delta}{\lambda}.
		\]
		Now compute the derivative of $\phi^*$:
		\[
		\left(\phi^*\right)'(x) = (g^*)'(x)-1 + \lambda\bigl(((g^*)^{-1})'(x)-1\bigr).
		\]
		Hence
		\[
		|(\phi^*)'(x)| \le |(g^*)'(x)-1| + \lambda|((g^*)^{-1})'(x)-1|
		\le \delta + \lambda\cdot\frac{\delta}{\lambda}=2\delta.
		\]
		Because $\int_{\T}\phi^*=0$ and $\phi^*$ is $1$-periodic, the Poincar\'e inequality gives
		\[
		\|\phi^*\|_{\infty}\le \frac12\|\left(\phi^*\right)'\|_{\infty}\le\delta.
		\]
		Therefore $\|\phi^*\|_{C^1}\le 2\delta+\delta=3\delta$, establishing (i).
		
		All the required properties are verified, with the constants
		$\alpha_1=\frac12$, $\alpha_2=c'/\lambda$, and the functions $\phi^*,\Psi^*$
		constructed above.
		
		
		
		
		
		
		\vspace{2em}
		
		
		
		
		\noindent\textbf{Data Availability Statement.}
		The authors state that this manuscript has no associated data and there is no conflict of interest.
		%	
		
		
		
		
		
		
		
		
		
		
		
		
		
		\begin{thebibliography}{DGOW}
			\renewcommand{\baselinestretch}{1.0}
			\setlength\itemsep{-1pt}
			%\def\bi#1{\bibitem[#1]{#1}}
			\small
			\bibitem[Arn11]{Ar11}  M.-C. Arnaud.
			{\it A non-differentiable essential irrational invariant curve
				for a $C^1$ symplectic twist map}. Journal of modern dynamics, 2011, 5 (3), pp.583-591.
			
			%\bibitem[Ban88]{B} V. Bangert. {\it Mather sets for twist maps and geodesics on tori}. Dynamics Reported {1} (1988), 1-45.
			
			\bibitem[BB13]{BB} P. Berger  and A. Bounemoura. {\it A geometrical proof of the persistence of normally hyperbolic submanifolds}. Dynamical Systems, 2013, 28 (4), 567-581.
			
			%\bibitem[BG99]{BG-1} A. Berretti, G. Gentile. {\it Scaling Properties for the Radius of Convergence of a Lindstedt Series: The Standard Map}. J. Math. Pures Appl. (9) 78 (1999), no. 2, 159-176.
			
			%\bibitem[BG00]{BG0} A. Berretti, G. Gentile. {\it Scaling properties for the radius of convergence of Lindstedt series: generalized standard maps}. J. Math. Pures Appl. (9) 79 (2000), no. 7, 691-713.
			
			%		\bibitem[BG01]{BG1} A. Berretti, G. Gentile. {\it Bryuno function and the standard map}. Comm. Math. Phys.,  220 (2001), no. 3, 623-656.
			%		
			%		\bibitem[BG02]{BG2} A. Berretti, G. Gentile. {\it Periodic and quasi-periodic orbits for the standard map}. Comm. Math. Phys. 231 (2002), no. 1, 135-156.
			
			%\bibitem[Bes01]{Bes} U. Bessi. {\it An analytic counterexample to the KAM theorem}. Ergodic Theory Dynam. Systems 20 (2000), 317-333.
			
			\bibitem[Bir20]{Bir20} G. D. Birkhoff.  {\it Surface transformations and their dynamical applications}. Acta Mathematicae 43 (1920) 44-47 Collected Mathematical Papers of G.D. Birkhoff , vol. II, p. 195-202, Dover, New York (1968) .
			
			%\bibitem[Bir32]{Bir32} G. D. Birkhoff. {\it Sur quelques courbes ferm\'{e}es remarquables}. Bull. Soc. Math. France 60 (1932), 1-26.
			
			
			
			\bibitem[BF19]{BF} A. Bounemoura and J. F\'{e}joz. {\it KAM, $\alpha$-Gevrey regularity and the $\alpha$-Bruno-R\"{u}ssmann condition}. Ann. Sc. Norm. Super. Pisa Cl. Sci. (5) 19 (2019), no. 4, 1225-1279.
			%
			%\bibitem[Bou19]{Bou} A. Bounemoura. {\it  Some remarks on the optimality of the Bruno-R\"{u}ssmann condition}. Bulletin de la soci\'{e}t\'{e}
			%math\'{e}matique de France, 2019, 147 (2), pp.341-353.
			%		
			%		
			\bibitem[CCD13]{CCD13} R. Calleja, A. Celletti and R. de la Llave. {\it A KAM theory for conformally symplectic systems: efficient algorithms and their validation}. J. Differ. Equ. {255} (2013), 978-1049.
			
			\bibitem[CCD20]{CCD20} R. Calleja, A. Celletti and R. de la Llave. {\it KAM Theory for Some Dissipative Systems}. New Frontiers of Celestial Mechanics: Theory and Applications. I-CELMECH 2020. Springer Proceedings in Mathematics \& Statistics, vol 399. Springer.
			
			
			\bibitem[CCD22]{CCD22} R. Calleja, A. Celletti and R. de la Llave. {\it KAM quasi-periodic solutions for the dissipative standard map}. Commun. Nonlinear Sci. Numer. Simul. 106 (2022), Paper No. 106111, 29 pp.
			%		
			%		
			%		%\bibitem[CdS]{CdS}
			%		%A. Cannas da Silva. {\it Lecture on Symplectic Geometry}. Lecture Notes in Mathematics, Vol. 1764, Springer-Verlag, Berlin (2001), xii+217 pp.
			%		
			%		%\bibitem[CF94]{CF} L. Chierchia and C. Falcolini. {\it A direct proof of a theorem by Kolmogorov in Hamiltonian systems}. Ann. Scuola Norm. Sup. Pisa Cl. Sci. (4) 21 (1994), no. 4, 541-593.
			%		
			%		
			\bibitem[Cas87]{Cas87} M. Casdagli. {\it Periodic orbits for dissipative twist maps}. Ergod. Th. \& Dynam. Sys. {7} (1987), 165-173.
			
			%\bibitem[CW13]{CW}C.-Q. Cheng and L. Wang. {\it Destruction of Lagrangian torus for positive definite Hamiltonian systems}. Geom. Funct. Anal. {23} (2013), 848-866.
			
			
			%\bibitem[Con78]{Co} C. Conley. {\it Isolated Invariant Sets and the Morse Index}. CBMS Regional Conference Series in Mathematics. {38} American Mathematical Society, Providence (1978).
			
			
			
			%\bibitem[dMvS93]{de}  W. de Melo, S. van Strien. {\it One-Dimensional Dynamics}. Springer-Verlag. (1993)
			
			%\bibitem[Den32]{De} A. Denjoy. {\it Sur les courbes d\'{e}finies par les \'{e}quations diff\'{e}rentielles \`{a} la surface du tore}. J. de Math. Pures et Appl., (9), 11 (1932), 333-375
			
			%\bibitem[Dev]{Dev}  R. Devaney. {\it An Introduction to Chaotic Dynamical Systems}. Wetsview Press. 2 (2003), 102-113
			
			%		\bibitem[Dav94]{Dav}  A. Davie. {\it The Critical Function for the Semistandard Map}. Nonlinearity 7,  (1994), 219-229.
			%
			%		%\bibitem[Eli86]{E1} H. Eliasson {\it Compensations of signs in a small divisor problem}. in Aspects dynamiques et topologiques des groupes infinis de transformation de la mecanique (Lyon, 1986), 37-48, Travaux en Cours Vol. 25, Hermann,
			%		% Paris, 1987.
			%		
			%		\bibitem[Eli96]{E2} H. Eliasson {\it Absolutely convergent series expansions for quasi-periodic motions}. University of Stockholm preprint (1988), and Math. Phys. Elect. J. 2, No. 4 (1996),
			%
			%		%\bibitem[Fa]{Fav} J. Favard. {\it Sur les meilleurs proc\'{e}d\'{e}s d'approximation}. B.S.M. 61 (1937), 243-256.
			%		
			%		%\bibitem[FK18]{FKO} B. Fayad and R. Krikorian. {\it   Some questions around quasi-periodic dynamics}. World Scientific Publishing Co. Pte. Ltd., Hackensack, NJ, 2018, 1909-1932.
			%		
			%		\bibitem[For94]{F} G. Forni. {\it Analytic destruction of invariant circles}. Ergod. Th. \& Dynam. Sys. {14} (1994), 267-298.
			%		
			%		
			%		
			%		%\bibitem[Gal93]{Ga1} G. Gallavotti. {\it Invariant tori: a field-theoretic point of view of Eliasson's work}. in Advances in dynamical systems and quantum physics (Capri, 1993), 117-132, World Sci. Publishing, River Edge, NJ, 1995.
			%		
			%		\bibitem[Gen15]{Ge1} G. Gentile. {\it Invariant curves for exact symplectic twist maps of the cylinder with Bryuno rotation numbers}. Nonlinearity 28 (2015) 2555–2585.
			
			%\bibitem[GM96]{GM} G. Gentile, V. Mastropietro. {\it Methods for the analysis of the Lindstedt series for KAM tori and renormalizability in classical mechanics}. A review with some applications, Rev. Math. Phys. 8 (1996), 393-444
			
			
			%  \bibitem[Gol01]{Gol01} C. Gol\'{e}. {\it Symplectic twist maps: global variational techniques}. Vol. 18. World Scientific, 2001.
			
			
			%\bibitem[Hal83]{Ha} G. Hall. {\it
				%Bifurcation of an attracting invariant circle: a Denjoy attractor}. Ergod.
			% Th. \& Dynam. Sys. {3} (1983), 87-118.
			
			
			
			%\bibitem[GLW25]{GLW} J. Guo, Y. Liu and L. Wang. {\it On non-KAM invariant circles for area-preserving twist maps}.  arXiv:2510.21098 [math.DS].
			
			
			
			\bibitem[Her79]{H79} M. R. Herman. {\it Sur la conjugation diff$\acute{e}$rentiable des diff$\acute{e}$omorphismes du cercle $\grave{a}$ des rotations}. Publ. Math. IHES {49} (1979), 5-233.
			
			
			\bibitem[Her83]{H83} M.R. Herman. {\it Sur les courbes invariantes par les diff\'{e}omorphismes de l'anneau}. Ast\'{e}risque {103-104} (1983), 1-221.
			%		
			%		\bibitem[Her86]{H33} M. R. Herman.  {\it Sur les courbes invariantes par les diff$\acute{e}$omorphismes de l'anneau}. Ast$\acute{\text{e}}$risque {144} (1986), 1-243.
			
			%\bibitem[H3]{H2} M.R. Herman. {\it In\'{e}galit\'{e}s ``a priori" pour des
				%tores lagrangiens invariants par des diff\'{e}omrphismes
				%symplectiques}. Inst. Hautes \'{E}tudes Sci. Publ. Math. {70}
			%(1990), 47-101.
			
			%\bibitem[Her90]{H3} M.R. Herman. {\it Non existence of Lagrangian graphs}.
			%unpublished preprint (1990).
			
			
			%\bibitem[Her98]{H98} M.R. Herman. {\it Some open problems in dynamical systems}. Proceedings of the International Congress of Mathematicians, Vol. II (Berlin, 1998) Doc. Math. 1998, Extra Vol. II, 797-808.
			
			
			\bibitem[HPS77]{HPS} M. Hirsch, C. Pugh and M. Shub. {\it Invariant manifolds}. Lecture Notes in Mathematics, Vol. 583. Springer-Verlag, (1977), ii+149 pp.
			
			%\bibitem[Jar29]{Jar29} V. Jarn\'{\i}k. Diophantischen Approximationen und Hausdorffsches mass. Mat. Sbornik 36 (1929), 371–382.
			
			
			%\bibitem[Khi64]{Khi64} A.~Ya.~Khinchin.{\it Continued Fractions}.University of Chicago Press, 1964.
			
			\bibitem[KH95]{KH} A. Katok and B. Hasselblatt. {\it Introduction to the Modern Theory of Dynamical Systems} (Encyclopedia of Mathematics and its Applications, 54). Cambridge University Press, Cambridge, 1995.
			%\bibitem[Laz93]{La} V. F. Lazutkin. {\it Recent results on the separatrix splitting for the standard map}. S\'{e}minaire de Th\'{e}orie spectrale et g\'{e}om\'{e}trie,  11(1992-1993), 157-165
			
			
			%\bibitem[LeC86]{Lec86} P.  Le Calvez. {\it Existence d'orbites quasi-periodiques dans les attracteurs de Birkhoff}. Commun. Math. Phys. {106} (1986), 383-394.
			
			\bibitem[LeC87]{Lec87} P.  Le Calvez. {\it Propri\'{e}t\'{e}s des attracteurs de Birkhoff}.
			Ergod. Th. \& Dynam. Sys. {8} (1987), 241-310.
			
			
			%\bibitem[Li]{Li} P. Libermann. {\it Sur les automorphismes infinit\'{e}simaux des structures symplectiques et des structures de contact}. (French) 1959 Colloque G\'{e}om. Diff. Globale
			%(Bruxelles, 1958) pp. 37-59 Centre Belge Rech. Math., Louvain.
			
			%\bibitem[LLW26]{LLW26}Q. Li, Y. Liu and L. Wang. {\it Persistence of  invariant graphs for twist maps under ultra-analytic perturbations}. Preprint.
			
			%\bibitem[LiW26]{LiW26}Q. Li and L. Wang. {\it On the Dynamics of Invariant Graphs for Dissipative Twist Maps}. Preprint.
			
			%\bibitem[LuW25]{LuW25} Y. Liu and L. Wang. {\it Herman's converse KAM mechanism revisited}. Preprint.
			
			%\bibitem[KO95]{KO}Y. Katznelson and D. Ornstein. {\it Smoothness of invariant curves}. Proceedings of the Conference in Honor of Jean-Pierre Kahane (Orsay, 1993) J. Fourier Anal. Appl.,Special Issue(1995), 283-310.
			
			
			
			\bibitem[Ma\~{n}78]{Man78} R. Ma\~{n}\'{e}. {\it Persistent manifolds are normally hyperbolic}. Trans. Amer. Math. Soc. 246 (1978), 261-283.
			
			%		\bibitem[MS03]{MS03} J.-P. Marco and D. Sauzin, Stability and instability for Gevrey quasi-convex near-integrable Hamiltonian systems, Publ. Math. IHES 96 (2003), 199–275.
			%		
			%		\bibitem[MS17]{MS17} S. Mar\`o and A. Sorrentino. {\it Aubry-Mather theory for conformally symplectic systems.} Comm. Math. Phys.  {354} (2017), 775-808.
			%		
			\bibitem[Mas23]{Mas23} J. E. Massetti. {\it Attractive invariant circles `a la Chenciner}. Regul. Chaot. Dyn. 28, 447-467 (2023).
			%
			%\bibitem[MF06]{MF06} J. Mather and G. Forni.{\it Action minimizing orbits in hamiltomian systems}. Transition to Chaos in Classical and Quantum Mechanics: Lectures given at the 3rd Session of the Centro Internazionale Matematico Estivo (CIME) held in Montecatini Terme, Italy, July 6–13, 1991. Berlin, Heidelberg: Springer Berlin Heidelberg, 2006. 92-186.
			%		
			%		\bibitem[Mo62]{Mo62} J. Moser. {\it On invariant curves of area-preserving mappings of an annulus}. Nachr. Akad. Wiss. G\"{o}ttingen,
			%		II. Math-Phys. KL. {1} (1962), 1-20.
			
			%\bibitem[Mo73]{Mo2}  J. Moser. {\it Stable and Random Motions in Dynamical Systems}. Princeton University
			%		Press. Princeton, N.J. (1973).
			
			%		\bibitem[MS03]{MS03} J. P. Marco and D. Sauzin. {\it Stability and instability for Gevrey quasi-convex near-integrable Hamiltonian systems}. Publ. Math. IHES {96} (2003), 199-275.
			
			%\bibitem[MY02]{MY} S. Marmi and J.-C. Yoccoz. {\it Some open problems related to small divisors} Lecture Notes in Math., 1784 Fond. CIME/CIME Found. Subser. Springer-Verlag, Berlin, 2002, 175-191.
			
			
			\bibitem[Mat84]{M84} J. N. Mather. {\it Non-existence of invariant circles}. Ergod. Th. \& Dynam. Sys. {4} (1984), 301-309.
			
			%\bibitem[Mat87]{M3} J. N. Mather. {\it Modulus of continuity for Peierls's barrier}. Periodic Solutions of Hamiltonian Systems and Related Topics. ed. P.H.Rabinowitz {\it et al}. NATO ASI Series C {209}. Reidel: Dordrecht, (1987), 177-202.
			
			%\bibitem[Mat88]{M4} J. N. Mather. {\it Destruction of invariant circles}. Ergod. Th. \& Dynam. Sys. {8} (1988), 199-214.
			
			
			%\bibitem[Mas23]{Ma} {J. Massetti. {\it Attractive invariant circles \`{a} la Chenciner}. Regul. Chaot. Dyn. 28, 447-467 (2023).}
			
			
			
			
			
			%\bibitem[Mos73]{Mo2}  J. Moser. {\it Stable and Random Motions in Dynamical Systems}. Princeton University
			%			Press. Princeton, N.J. (1973).
			%\bibitem[New04]{Ne} S. Newhouse. {\it Cone-fields, domination, and hyperbolicity}. Brin,
			% Michael (ed.) et al., Modern dynamical systems and applications.
			% Dedicated to Anatole Katok on his 60th birthday. Cambridge University Press. 419-432, 2004.
			\bibitem[New04]{Ne} S. Newhouse. {\it Cone-fields, domination, and hyperbolicity}. Brin,
			Michael (ed.) et al., Modern dynamical systems and applications.
			Dedicated to Anatole Katok on his 60th birthday. Cambridge University Press. 419-432, 2004.		
			
			
			%\bibitem[OP08]{OP} A. Olvera and N. Petrov. {\it Regularity Properties of Critical Invariant Circles
				%of Twist Maps, and Their Universality}. SIAM J. Appl. Dyn. Syst.   7 (2008), 962-987.
			
			
			
			
			%\bibitem[P]{P} J. P\"{o}schel.  {\it \"{U}ber invariante Tori in differenzierbaren Hamiltonschen Systemen}. (German)[On invariant tori in differentiable Hamiltonian systems]
			%Bonner Math. Schriften, 120[Bonn Mathematical Publications]
			%Universit\"{a}t Bonn, Mathematisches Institut, Bonn, 1980. 103 pp.
			
			
			%\bibitem[R\"{u}s70]{R70} H. R\"{u}ssmann.  {\it Kleine Nenner. I. \"{U}ber invariante Kurven differenzierbarer Abbildungen eines Kreisringes}. Nachr. Akad. Wiss. G\"{o}ttingen Math.-Phys. Kl. II(1970), 67-105.
			
			
			
			%\bibitem{R1} H. R\"{u}ssman.  {\it On optimal estimates for the solutions of linear partial differential equations of first order with constantm coefficients on the torus}. Lecture Notes in Physics. {38} (1975), 598-624.
			
			%		\bibitem[R\"{u}s83]{R2} H. R\"{u}ssmann. {\it On the existence of invariant curves of twist mappings of an annulus}. Lecture Notes in Mathematics.  {1007} (1983),  677-718.
			
			
			%		\bibitem[R\"{u}s91]{R3} H. R\"{u}ssman. {\it On the frequencies of quasi-periodic solutions of nearly integrable Hamiltonian systems}. Euler International Mathematical Institute, St. Petersburg, Dynamical Systems, 14-27 October 1991.
			
			%\bibitem[Sal04]{Sa} D. Salamon. {\it The Kolmogorov-Arnold-Moser theorem}. Math. Phys. Electron. J. 10 (2004), Paper 3, 37 pp.		
			
			%\bibitem[Sie42]{Sie} C. Siegel. {\it Iteration of analytic functions}. Ann. of Math. (2) 43 (1942), 607-612.
			
			
			\bibitem[SW26]{SoW} A. Sorrentino and L. Wang. {\it On the Destruction of Invariant Lagrangian Graphs for Conformal Symplectic Twist Maps}.  Calc. Var. Partial Differential Equations,
			65(5), 2026, Article 153.
			
			%\bibitem[SuW]{SW} X. Su and L. Wang. {\it Total destruction of invariant tori for the generalized Frenkel-Kontorova model}, {Journal of Math. Phys.}, 53 (2012), 022702, 7pp.
			
			%\bibitem[Ta]{T} F. Takens. {\it A $C^1$ counterexample to Moser's twist theorem. } Nederl. Akad. Wetensch. Proc. Ser. A {74}. Indag. Math.  {33}
			%(1971), 378-386.
			
			%\bibitem[Wan12]{W} L. Wang. {\it Variational destruction of invariant circles}. {Discrete Contin. Dyn. Syst.}, 32 (2012), 4429-4443.
			
			
			%\bibitem[Wan15]{W15} L. Wang. {\it Destruction of invariant circles for Gevrey area-preserving twist maps}. {J. Dynam. Differential Equations},  27 (2015), 283-295.
			
			%\bibitem[W2]{W2}L. Wang. {\it Total destruction of Lagrangian tori}. J. Math. Anal. Appl. {410} (2014), 827-836.012).
			
			
			%\bibitem[Wan22]{W22} L. Wang. {\it Quantitative destruction of invariant circles}. {Discrete Contin. Dyn. Syst.},  42 (2022), 1569-1583.
			
			
			%\bibitem[Wan23]{W23} L. Wang. {\it Quantitative Destruction and Persistence of Lagrangian Torus in Hamiltonian Systems}. arXiv:2312.01695v3 [math.DS]
			
			
			%\bibitem[Yoc84]{Y1} J.-C. Yoccoz. {\it Conjugaison diff\'{e}rentiable des diff\'{e}omorphismes du cercle dont le nombre de rotation v\'{e}rifie une condition diophantienne}. Ann. Sci. \'{E}cole Norm. Sup. (4) 17 (1984), no. 3, 333-359.
			
			%\bibitem[Yoc92]{Y2} J.-C. Yoccoz. {\it Travaux de Herman sur les tores invariants. }Astérisque, 206, 311–346.
			
		\end{thebibliography}
		
		
		
		
		
		
		
		
	\end{document}
	
	
}

